\documentclass[10pt,a4paper]{article}

% ----------------- Paquetes -------------------%
\usepackage[T1]{fontenc}
\usepackage[utf8]{inputenc}
\usepackage[a4paper,margin=2.5cm]{geometry}
\usepackage{mathptmx}             
\usepackage{changepage} 
\usepackage{setspace}
\usepackage[spanish]{babel}
\usepackage[labelfont=bf,labelsep=space]{caption}
\usepackage{amssymb}
\usepackage{amsmath}
\usepackage{ebgaramond}
\usepackage{placeins}
\usepackage{etoolbox}
\usepackage{setspace}
\usepackage{parskip} 
\usepackage{tcolorbox}
\usepackage{needspace}
\usepackage{xparse}
\usepackage{ocgx2}
\usepackage{hyperref}
\usepackage{hypcap}
\usepackage{multirow}
\usepackage{soul}\sethlcolor{yellow}% resaltado amarillo: \hl{texto} (Panel TCEE, ⌃H)


%%%%%%%%%%%%%%%%%%%%%%%%%%%%%%%%%%%%%%%%%%%%%%%%%%%%%%%%%%%%%%%%
% --- Fija primero tus valores globales "buenos" ---
\setstretch{1.2}                 % Interlineado global
\setlength{\parskip}{0.7\baselineskip}
\setlength{\parindent}{0pt}
\newlength{\GlobalParskip}
\setlength{\GlobalParskip}{\parskip}

\newcounter{eqnum}[subsection]
\renewcommand{\theeqnum}{\arabic{eqnum}}
\newcounter{alphaitem}
\newcounter{numitem}

%% ------------------- Recuadros----------------------%%
\newcommand{\puntosclave}[1]{%
  \begin{tcolorbox}[
    colframe=black,
    title={Puntos clave},
    before upper={\setlength{\parindent}{0.5cm}\setlength{\parskip}{0.5\baselineskip}}
  ]
  #1
  \end{tcolorbox}
}

\newcommand{\modificaciones}[1]{
  \begin{tcolorbox}[
    colback=white,
    colframe=red,
    title={Modificaciones a realizar},
    before upper={\setlength{\parindent}{0.5cm}\setlength{\parskip}{0.5\baselineskip}}
  ]
  #1
  \end{tcolorbox}
}

%% ------------------- CITA ----------------------%%
% \begin{cita}[Autor][Año][Obra] texto \end{cita}: cita sangrada y, debajo a la derecha, «AUTOR (Año), Obra». Los tres datos son opcionales.
% En el Panel TCEE: escribir \cita e Intro (el cursor va al texto; Tab pasa a Autor, Año y Obra).
\NewDocumentEnvironment{cita}{ O{} O{} O{} +b }{%
  \begin{adjustwidth}{2cm}{0pt}
  \raggedright
  #4\par
  \raggedleft
  \if\relax\detokenize{#1#2#3}\relax
    % sin firma
  \else
    #1%
    \if\relax\detokenize{#2}\relax\else\if\relax\detokenize{#1}\relax\else\space\fi(#2)\fi
    \if\relax\detokenize{#3}\relax\else\if\relax\detokenize{#1#2}\relax\else, \fi\textit{#3}\fi
  \fi
  \end{adjustwidth}
}{}



%% ------------------- Listas----------------------%%
    % --- Lista alfabética ---
    \newenvironment{la}
      {\begin{list}{\alph{alphaitem})}{%
          \usecounter{alphaitem}%
          \setlength{\leftmargin}{1cm}%
          \setlength{\parskip}{3pt}% 
          \setlength{\itemsep}{0pt}% ← espacio entre ítems
          \setlength{\parsep}{0pt}%  ← párrafos dentro de un ítem
      }}
      {\end{list}}
      
    % --- Lista numérica ---
    \newenvironment{lnum}
      {\begin{list}{\arabic{numitem}.}{%
          \usecounter{numitem}%
          \setlength{\leftmargin}{1cm}%
          \setlength{\parskip}{3pt}% 
          \setlength{\itemsep}{0pt}% ← espacio entre ítems
          \setlength{\parsep}{0pt}%  ← párrafos dentro de un ítem
      }}
      {\end{list}}
      
% ------------------- Imágenes----------------------%
    \graphicspath{{Img/}}% Rutas por defecto 
    % --- Imagen adaptativa ---
    % Registros de dimensión definidos una sola vez
    \newdimen\imgcap
    \newdimen\imgmin
    \newdimen\imgmax
    \newdimen\imgremain
    \newdimen\imgh
    \newdimen\imgneed
    
    \newcommand{\imagenfit}[3]{%
      \addvspace{10pt}% separación antes
      \par\begingroup
        % Parámetros de composición (asignaciones, no \newdimen)
        \imgcap=1\baselineskip            % alto estimado de título+fuente
        \imgmin=0.15\textheight
        \imgmax=0.15\textheight
        \imgremain=\pagegoal
        \ifdim\pagegoal=\maxdimen \imgremain=0pt\fi
        \advance\imgremain by -\pagetotal
        \advance\imgremain by -\imgcap
        % Decisión de altura destino
        \ifdim\imgremain<\imgmin
          \newpage
          \imgh=0.20\textheight
        \else
          \imgh=\imgremain
          \ifdim\imgh>\imgmax \imgh=\imgmax \fi
        \fi
        % Reservar espacio para evitar cortes del bloque completo
        \imgneed=\imgh \advance\imgneed by \imgcap
        \Needspace{\imgneed}
        % Cuerpo
        \noindent\begin{minipage}{\textwidth}
          \centering
          \captionof{figure}{\textemdash\ #2}\par\vspace{0.3em}% título arriba
          \includegraphics[width=\linewidth,height=\imgh,keepaspectratio]{\detokenize{#1}}\par
          \vspace{0.3em}\small\textit{Fuente: }#3% fuente abajo
        \end{minipage}%
      \endgroup
      \par\addvspace{10pt}%
    }

%------------ Bloque de RECUADRO ESPECIAL -------------%
    \newenvironment{specialblock}[1]{%
      \begin{quote} % crea sangría a izquierda y derecha
      \small % texto más pequeño
      \noindent\textbf{#1}\\[-0.3em] % título en negrita
      \rule{\linewidth}{0.4pt}\\[0.6em]}{
      \end{quote}}
      

%-------------------------- Author Font -----------------------%
    \newcommand{\authorfont}[1]{{\fontfamily{EBGaramond-TLF}\selectfont #1}}

%-------------------------- Anexos -----------------------%
    \makeatletter
    \newcounter{anexo}
    \renewcommand{\theanexo}{\Roman{anexo}}
    \newcommand{\anexo}[1]{%
      \clearpage
      \phantomsection
      \refstepcounter{anexo}%
      \noindent\textbf{Anexo \theanexo.- }#1\par\medskip
      \addcontentsline{stc}{section}{Anexo \theanexo.- #1}%
    }
    \makeatother

    %-------------------------- Lista de figuras (personal) -----------------------%
\makeatletter
\newcommand{\MyListOfFigures}{%
  \clearpage
  \phantomsection
  {\centering\bfseries\Large Índice de figuras\par}%
  \medskip
  % Entrada en nuestro índice de contenido (stc)
  \addcontentsline{stc}{section}{Índice de figuras}%
  % Recuperación del listado (sin cabecera propia)
  \begingroup
    \parindent=0pt\parskip=2pt
    \@starttoc{lof}%
  \endgroup
}
\makeatother

%-------------------------- Bibliografía -----------------------%
\makeatletter
% Contador para las referencias
\newcounter{myref}
% Cabecera de la bibliografía, entra en el índice 'stc'
\newcommand{\MyBibliography}{%
  \clearpage
  \phantomsection
  {\centering\bfseries\Large Bibliografía y referencias\par}%
  \medskip
  \addcontentsline{stc}{section}{Bibliografía y referencias}%
  \setcounter{myref}{0}%
}
\newcommand{\MyBibItem}[4]{%
  \refstepcounter{myref}%
  \noindent[\themyref]\ #1.\ \emph{#2}. #3. #4\par
}
\makeatother

%--------- Establecer cuatro niveles de índice --------%
    \setcounter{tocdepth}{4}   
    \setcounter{secnumdepth}{4}% numera hasta \paragraph

%%%%%%%%%%%%%%%%%%%%%%%%%%%%%%%%

\makeatletter

    %--------- Interlineado Global --------%
    \edef\GlobalBaseLineStretch{\baselinestretch}
    
    %--------- Espaciado de seccinones --------%
    \renewcommand\section{\@startsection{section}
      {1}{\z@}
      {2.5ex \@plus 1ex \@minus .2ex} % espacio antes
      {1.5ex \@plus .2ex}             % espacio después
      {\normalfont\Large\bfseries}    % formato del título
    }
    \renewcommand\subsection{\@startsection{subsection}{2}{\z@}
      {3ex \@plus 0.8ex \@minus .2ex}
      {1.5ex \@plus .2ex}
      {\normalfont\large\bfseries}
    }
    \renewcommand\subsubsection{\@startsection{subsubsection}{3}{\z@}
      {3ex \@plus 0.5ex \@minus .2ex}
      {1ex \@plus .2ex}
      {\normalfont\normalsize\bfseries}
    }
    \renewcommand\paragraph{\@startsection{paragraph}{4}{\z@}%
    {-2ex \@plus -0.5ex \@minus -.2ex}% espacio antes
    {0.8ex \@plus .2ex}% espacio después
    {\normalfont\normalsize\itshape}}% ← cursiva en lugar de negrita

    %--------- Ecuaciones --------%   
    % Variables globales para almacenar títulos y notas
    \gdef\leq@current@section{}%
    \gdef\leq@current@subsection{}%
    \gdef\leq@last@printed@section{}%
    \gdef\leq@last@printed@subsection{}%
    
    % Comando para añadir nota (invisible en el cuerpo)
    \newcommand{\eqnote}[1]{%
      \addcontentsline{leq}{leqnote}{#1}%
    }
    
    % Comando para ecuaciones
    \newcommand{\eqblock}[2]{%
      % Verificar si necesitamos imprimir el título de sección
      \ifx\leq@current@section\leq@last@printed@section\else
        \addcontentsline{leq}{leqsection}{\leq@current@section}%
        \global\let\leq@last@printed@section\leq@current@section
        \gdef\leq@last@printed@subsection{}% Reset subsección al cambiar sección
      \fi
      % Verificar si necesitamos imprimir el título de subsección
      \ifx\leq@current@subsection\leq@last@printed@subsection\else
        \ifx\leq@current@subsection\@empty\else
          \addcontentsline{leq}{leqsubsection}{\leq@current@subsection}%
          \global\let\leq@last@printed@subsection\leq@current@subsection
        \fi
      \fi
      % Ecuación
      \refstepcounter{eqnum}%
      \begin{equation*}
        #1\tag{E.\theeqnum}
      \end{equation*}%
      \addcontentsline{leq}{leqt}{[E.\theeqnum] -- #2}%
      \addcontentsline{leq}{leqm}{#1}%
    }
    
    %%% ----- Índice de ecuaciones ----------%%%
    % Tamaño para []
    \newcommand{\leqIndexFont}{\normalsize}
    \newcommand{\leqTitleFont}{\tiny}
    \newcommand{\leqNoteFont}{\tiny}
    % Cabecera de sección 
    \newcommand*\l@leqsection[2]{%
      \addvspace{25pt}% Mayor separación antes de nueva sección
      \noindent\leqIndexFont
      \makebox[\linewidth]{\rule{.38\linewidth}{0.4pt}\ \ \textbf{  \MakeUppercase{#1}}\ \ \rule{.38\linewidth}{0.4pt}}%
      \par\vspace{-10pt}
    }
    % Cabecera de subsección 
    \newcommand*\l@leqsubsection[2]{%
      \par\addvspace{15pt}%
      \noindent\leqIndexFont #1
      \par\vspace{-7pt}
    }
    % Título de ecuación 
    \newcommand*\l@leqt[2]{%
    \par\addvspace{10pt}%
      \par\noindent\leqTitleFont #1\par
    }
    % Miniatura de ecuación
    \newcommand*\l@leqm[2]{%
      \vspace{0.3\baselineskip}%
      \par\scriptsize\noindent\hspace*{1.5em}\(#1\)\par%
    }
    % Nota de ecuación 
    \newcommand*\l@leqnote[2]{%
      \vspace{0.2\baselineskip}%
      \par\noindent\leqNoteFont\hspace*{1.5em}\textbf{Nota.--} #1\par\vspace{0.5\baselineskip}%
    }
    % Generación del índice
    \newcommand{\mylistofequations}{%
      \clearpage
      \phantomsection
      % Título visual de la lista (ajústalo a tu gusto):
      {\centering\bfseries\Large Índice de ecuaciones\par}
      % Entrada en nuestro índice 'stc':
      \addcontentsline{stc}{section}{Índice de ecuaciones}%
      % Llamada a tu lista real (usa el nombre exacto de tu macro):
      \@starttoc{leq}%
    }
    
    %--------- Captura de títulos de sección --------%
    \let\leq@original@section\section
    \let\leq@original@subsection\subsection
    % Flag para saber si estamos en una sección asterisco
    \newif\ifLeqStarSection
    \LeqStarSectionfalse
    
    % Redefinir \section CON CONTROL DE ASTERISCO
    \renewcommand{\section}{%
      \@ifstar{\leq@section@star}{\@dblarg\leq@section@nostar}%
    }
    \def\leq@section@star#1{%
      \global\LeqStarSectiontrue
      \gdef\leq@current@section{#1}%
      \gdef\leq@current@subsection{}%
      \leq@original@section*{#1}%
      \global\LeqStarSectionfalse
    }
    \def\leq@section@nostar[#1]#2{%
      \global\LeqStarSectionfalse
      \gdef\leq@current@section{#2}%
      \gdef\leq@current@subsection{}%
      \leq@original@section[#1]{#2}%
    }
    % Redefinir \subsection
    \renewcommand{\subsection}{%
      \@ifstar{\leq@subsection@star}{\@dblarg\leq@subsection@nostar}%
    }
    \def\leq@subsection@star#1{%
      \global\LeqStarSectiontrue
      \gdef\leq@current@subsection{#1}%
      \leq@original@subsection*{#1}%
      \global\LeqStarSectionfalse
    }
    \def\leq@subsection@nostar[#1]#2{%
      \global\LeqStarSectionfalse
      \gdef\leq@current@subsection{#2}%
      \leq@original@subsection[#1]{#2}%
    }

    % ----- Restauración de valores Globales de estilo ----------%
    % Flags
    \newif\ifInFootnote
    \InFootnotefalse
    \pretocmd{\@footnotetext}{\InFootnotetrue}{}{}
    \apptocmd{\@footnotetext}{\InFootnotefalse}{}{}
    % Profundidad de listas (soporta anidadas)
    \newcount\MyListDepth \MyListDepth=0
    \AtBeginEnvironment{la}{\advance\MyListDepth by 1}
    \AfterEndEnvironment{la}{\advance\MyListDepth by -1}
    \AtBeginEnvironment{lnum}{\advance\MyListDepth by 1}
    \AfterEndEnvironment{lnum}{\advance\MyListDepth by -1}
    \AtBeginEnvironment{itemize}{\advance\MyListDepth by 1}
    \AfterEndEnvironment{itemize}{\advance\MyListDepth by -1}
    \AtBeginEnvironment{enumerate}{\advance\MyListDepth by 1}
    \AfterEndEnvironment{enumerate}{\advance\MyListDepth by -1}
    % Restauración diferida: hazlo al INICIO del próximo párrafo real
    \newif\if@pendingRestore
    \@pendingRestorefalse
    \newcommand{\RequestRestore}{\global\@pendingRestoretrue}
    \everypar{%
      \if@pendingRestore
        \global\@pendingRestorefalse
        \ifInFootnote\else
          \ifnum\MyListDepth=0
            \setlength{\parskip}{\GlobalParskip}%
            \setstretch{\GlobalBaseLineStretch}%
          \fi
        \fi
      \fi
    }
    % Al final de bloques:
    \AfterEndEnvironment{equation}{\RequestRestore}
    \AfterEndEnvironment{equation*}{\RequestRestore}
    \AfterEndEnvironment{la}{\RequestRestore}
    \AfterEndEnvironment{lnum}{\RequestRestore}
    % Al final de macros:
    \apptocmd{\eqblock}{\RequestRestore}{}{}
    \apptocmd{\imagenfit}{\RequestRestore}{}{}
    
\makeatother

% ===================== ToC propio con 4 niveles (único bloque) =====================
\makeatletter

% --- Escritores: añadimos una línea al .stc en cada cabecera numerada ---
% Guardamos las versiones actuales para llamar al comportamiento normal
\let\MyTOC@orig@section\section
\let\MyTOC@orig@subsection\subsection
\let\MyTOC@orig@subsubsection\subsubsection
\let\MyTOC@orig@paragraph\paragraph

% SECTION
\renewcommand{\section}{%
  \@ifstar{\MyTOC@section@star}{\@dblarg\MyTOC@section@nostar}%
}
\def\MyTOC@section@star#1{%
  \MyTOC@orig@section*{#1}% sin entrada al .stc
}
\def\MyTOC@section@nostar[#1]#2{%
  \MyTOC@orig@section[{#1}]{#2}% comportamiento normal (escribe al .toc)
  % Escribimos también nuestra línea al .stc (texto con \numberline)
  \addcontentsline{stc}{section}{\protect\numberline{\thesection}#2}%
}

% SUBSECTION
\renewcommand{\subsection}{%
  \@ifstar{\MyTOC@subsection@star}{\@dblarg\MyTOC@subsection@nostar}%
}
\def\MyTOC@subsection@star#1{%
  \MyTOC@orig@subsection*{#1}%
}
\def\MyTOC@subsection@nostar[#1]#2{%
  \MyTOC@orig@subsection[{#1}]{#2}%
  \addcontentsline{stc}{subsection}{\protect\numberline{\thesubsection}#2}%
}

% SUBSUBSECTION
\renewcommand{\subsubsection}{%
  \@ifstar{\MyTOC@subsubsection@star}{\@dblarg\MyTOC@subsubsection@nostar}%
}
\def\MyTOC@subsubsection@star#1{%
  \MyTOC@orig@subsubsection*{#1}%
}
\def\MyTOC@subsubsection@nostar[#1]#2{%
  \MyTOC@orig@subsubsection[{#1}]{#2}%
  \addcontentsline{stc}{subsubsection}{\protect\numberline{\thesubsubsection}#2}%
}

% PARAGRAPH
\renewcommand{\paragraph}{%
  \@ifstar{\MyTOC@paragraph@star}{\@dblarg\MyTOC@paragraph@nostar}%
}
\def\MyTOC@paragraph@star#1{%
  \MyTOC@orig@paragraph*{#1}%
}
\def\MyTOC@paragraph@nostar[#1]#2{%
  \MyTOC@orig@paragraph[{#1}]{#2}%
  % Si no numeras los \paragraph, cambia \theparagraph por texto plano sin \numberline
  \addcontentsline{stc}{paragraph}{\protect\numberline{\theparagraph}#2}%
}

% --- Impresor del índice propio (lee el .stc) ---
\newcommand{\MyTOC}{%
  \par\bigskip
  {\centering\bfseries\Large Índice de contenido\par}%
  \medskip
  \begingroup
    \parindent=0pt\parskip=2pt
    \@starttoc{stc}%
  \endgroup
}

\makeatother
% ===================== fin ToC propio =====================


%%%%%%%%%%%%%%


% --- Datos del tema ---
\title{Teoría impositiva (V): Imposición óptima.\\
\large Tipo impositivo óptimo. Regla de Ramsey. El compromiso entre eficiencia y equidad de la Política Tributaria}
\author{Marcelo Ortega}
\date{Agosto 2026}

\begin{document}
\maketitle
\newpage

% --- Página de resumen y modificaciones ---
\puntosclave{ %% Escribir puntos clave Apuntes CECO
     
    }
\vspace{1cm}

\modificaciones{
    
    }

\newpage
\MyTOC
\newpage

\mylistofequations
\clearpage

\MyListOfFigures
\clearpage


\pagenumbering{arabic}
\setcounter{page}{1}


\section{Introducción}



\subsection{Contextualización}


\subsubsection{Relevancia}




\subsubsection{Problemática}



\subsection{Estructura de la exposición}




\clearpage
\section{Teoría de la Imposición Óptima: Marco conceptual}
\puntosclave{
    
}

HARBERGER, uno de los pioneros del modelo de incidencia fiscal escribió una vez sobre su experiencia en Indonesia, donde los automóviles están gravados más fuertemente que las motocicletas.

Esta diferencia impositiva proporcionaba un gran incentivo para hacer que las motocicletas se parecieran más a los automóviles. Como informa HARBERGER: \textit{los ciclos de tres ruedas eran convertidos, mediante ingeniosos añadidos, en autobuses virtuales, o al menos en taxis. A veces se añadía un solo banco, con el pasajero mirando hacia atrás. Otras veces el ciclo era alargado por la parte trasera, con dos bancos que recorrían cada lado, y quizá incluso con un pequeño estribo adicional que atravesaba lateralmente la parte posterior (donde estaría el parachoques trasero del automóvil). Debo decir que me quedé verdaderamente asombrado cuando vi mi primera motocicleta para ocho pasajeros}.

Este ejemplo pone de relieve un hecho sencillo: los mercados no aceptan los impuestos sin reaccionar. Si existe alguna acción que los participantes del mercado puedan emprender para minimizar la carga de los impuestos, la emprenderán. Mientras haya sustitutos para el consumo de cualquier bien gravado, algunos consumidores se desplazarán hacia esos sustitutos para evitar el impuesto, y mientras haya alternativas a la producción de bienes gravados, algunos productores pasarán a producir esas alternativas.

En temas anteriores, [\authorfont{Ver Tema 4.B.9}], aprendimos cómo los intentos de minimizar las cargas fiscales tienen costes de eficiencia para la sociedad. En ausencia de fallos de mercado, la eficiencia social se maximiza en el equilibrio competitivo sin intervención gubernamental. Cuando el gobierno grava a los participantes del mercado, estos cambian su comportamiento para evitar el impuesto y alejan el mercado del equilibrio competitivo, reduciendo así la eficiencia social. Dicho sencillamente, para la sociedad es costoso transportar a personas en peligrosas motocicletas para ocho pasajeros en lugar de en automóviles.

No obstante, teniendo en cuenta que cualquier gobierno necesita recaudar un cantidad suficiente que le permita llevar a cabo las labores propias que la sociedad le ha otorgado, lo que procede preguntarnos es si existe alguna forma de recaudar estos impuestos minimizando las distorsiones generadas por la intervención. Para abordar esta cuestión, presentaremos los principales resultados de la Teoría de la Imposición Óptima. 






 
La Teoría de fa Imposición Óptima (TIO) constituye uno de los temas clásicos de cualquier programa o manual de Economía de la Imposición. Dentro de esta. no resulta infrecuente encontrar los contenidos p1incipales de la TIO integrados con los dedicados al principio de neutralidad. las distorsiones impositivas y la medición del exceso de gravamen. La explicación está en que. en su concepción inicial, la TIO abordaba esencialmente el diseño óptimo de la imposición sobre el consumo. sin considerar ninguna meta distributiva. por lo que las reglas propuestas tenían como único objetivo la minimización del exceso de gravamen. Tras el desarrollo experimentado por la TIO en la década de los años setenta del siglo XX. con la incorporación de modelos de imposición óptima con metas redistributivas (a partir de los trabajos seminales de MiITlees y de Atkinson y Stiglitz), la TlO pasa a diferenciar dos grandes bloques, el dedicado a las reglas de imposición óptima sobre el consumo (sin metas redisributivas) y el dedicado a las reglas de la imposición óptima sobre la renta personal, extendido al gasto personal como una alternativa al gravamen de la renta.

En la confección del presente tema hemos optado por incorporar, tras la Introducción. un segundo apa11ado dedicado a presentar, desde un enfoque analítico. el principio de neutralidad. entroncado en el marco teórico de la Economía del Bienestar. y en concreto, de sus dos Teorema Fundamentales. A pa11ir de la noción de distorsión impositiva, en este apartado se incluye una exposición sintética del concepto de exceso de gravamen. esencial para la cuantificación monetaria de los costes de eficiencia generados por la imposición distorsionante y, consecuentemente, para la comprensión de los modelos que confo1man la TIO. La imp01tancia actual en la Economía Pública y, especialmente, en la Economía de la Imposición del concepto del Coste Marginal de los Fondos Públicos. ha recomendado su incorporación al final de este segundo apartado. Los apartados tres y cuatro siguen la estructura habitual, antes comentada, de separar las reglas de imposición óptima según estén dirigidas a la imposición sobre el consumo, en ausencia de metas distributivas (apartado 3), o incorporen un o��jetivo redistributivo, con el impuesto sobre la renta personal como eje principal de su formulación.

Dentro de que estamos ante un tema. como hemos dicho, clásico, el apartado 5 ha tratado de incorporar los contenidos que han pasado más recientemente a considerarse, con generalidad, como pa1te de la TIO, sí bien en algunos casos se trata de contenidos cuyas aportaciones seminales tienen cinco décadas. Esto se explica por la necesidad de incorporar el principio de sencillez de gestión a la formulación de reglas de diseño impositivo óptimo. con el propósito de acercar la TlO hacia escenarios de reforma fiscal más realistas, donde la aplicación de los impuestos por parte de la administración tributaria y la existencia de comportamientos de evasión y elusión fiscal por los· contribuyentes son elementos que no se puede obviar. En este sentido, se ha hecho un esfuerzo de síntesis, en la medida que las aportaciones científicas en este terreno son abundantes, si bien hemos optado por ofrecer los resultados básicos más generales. El tema concluye, precisamente con una valoración sobre el alcance que hoy en día puede tener la TIO en losprocesos de reforma fiscal abordados por los gobiernos.


\subsection{Dimensiones de la imposición pública}
La Economía de la Imposición es una disciplina fundamental para comprender cómo los impuestos afectan a la economía, desde múltiples perspectivas. y cómo esto debería tenerse en cuenta a la hora de diseñar los sistemas tributarios. La forma más habitual de abordar el estudio de esta materia es organizar los diferentes temas tratados poniéndolos en relación .con los denominados principios impositivos, encargados de orientar el diseño y la aplicación de sistemas impositivos neutrales, en el sentido de eficientes, equitativos, sencillos de administrar, capaces de adaptarse con flexibilidad a los ciclos económicos y, finalmente, suficientes para financiar las necesidades de gasto público demandadas por la sociedad. Puede decirse que la economía de la imposición proporciona principios fundamentales para guiar el diseño y la aplicación de sistemas impositivos adecuados a la realidad económica y social de los países.

\subsubsection{Dimensión distributiva: Equidad}


\subsubsection{Dimesión asignativa: Eficiencia}
En la valoración de los sistemas impositivos desde el punto de vista de la eficiencia económica, tres conceptos emergen como piedras angulares: el principio de neutralidad, las distorsiones impositivas y el exceso de gravamen. Estos conceptos, intrínsecamente relacionados, ofrecen una lente a través de la cual podemos analizar y evaluar el impacto de los sistemas tributarios en la eficiencia económica y la equidad.

El principio de neutralidad, en esencia, busca que los impuestos no interfieran con las decisiones económicas naturales de los individuos y las empresas. En un mundo ideal. los impuestos deberían ser diseñados de manera que los agentes económicos tomen decisiones basadas únicamente en las consideraciones económicas. sin verse influidos por consideraciones fiscales. Sin embargo, la realidad es que los impuestos introducen distorsiones en el mercado que pueden alterar las elecciones de consumo, ahorro, trabajo, inversión y producción, incluso otras decisiones más personales. como las de casarse, tener hijos, o dónde residir.

La Teoría de la Imposición Óptima (TIO) constituye un marco analítico cuyo propósito es encontrar las reglas que deberían orientar el diseño impositivo para. de acuerdo con los principios de la imposición, maximizar el bienestar social. En concreto. el fin principal de la TIO es determinar cómo debería ser la estructura de los impuestos -fundamentalmente, los tipos de gravamen a emplear y las magnitudes a gravar- con los que un gobierno que aspira a maximizar el bienestar social debería recaudar los ingresos fiscales necesarios para financiar los gastos públicos. minimizando los costes de eficiencia causados por las distorsiones económicas generadas y atendiendo, en su caso, a las metas redistributivas derivadas de las preferencias sociales.

Inicialmente, la TIO se centró, a partir del trabajo seminal de RAMSEY (1927). en estudiar las reglas que debían cumplir los impuestos sobre el consumo bajo la consideración de que el principio de neutralidad (eficiencia) era el único que debía inspirar su diseño óptimo. además. por supuesto, del de suficiencia, que operaba a modo de restricción recaudatoria de los problemas de optimización formulados. La ausencia de metas re distributivas suponía considerar un gobierno guiado por unas preferencias sociales de tipo utilitarista o lineales (benthwnistas). Tras sucesivas aportaciones para extender las reglas de imposición óptima sobre el consumo, la TlO no incorporó en sus modelos el dilema entre eficiencia y equidad hasta el trabajo seminal de MIRRLEES (1971), lo que permitió hablar de una segunda ola de modelos de imposición óptima. En un escenario de información imperfecta, MIRRLEES formalizó el problema del planificador social que se enfrenta a la heterogeneidad no observada que presentan los contribuyentes en relación con su capacidad innata para obtener ingresos. En su modelo original, el gobierno maximizador del bienestar social puede observar los ingresos de los ciudadanos, los cuales dependen, en cada caso. tanto de la capacidad innata como del esfuerzo, pero sin poder separar la contribución de ambos. Así, se aspira a gravar a los individuos de alta capacidad, estos podrán verse desincentivados a esforzarse para ganar esos ingresos. Al considerar esa heterogeneidad no observada, en un escenario con utilidad marginal decreciente del consumo, la fijación de los tipos óptimos debe tener en cuenta los efectos de desincentivo, lo que sitúa el diseño impositivo frente al conocido dilema entre igualdad y eficiencia. A partir de esta contribución seminal, la TIO ha desarrollado múltiples extensiones. las cuáles han ido incorporando nuevos elementos del diseño impositivo con el propósito de acerca el marco analítico de la TIO a la realidad de las figuras incluidas en los sistemas fiscales aplicados en el mundo real. En esta evolución, ha destacado la inclusión en los marcos analíticos del resto de principios impositivos. especialmente el principio de sencillez de gestión, incorporando tanto el cumplimiento de los contribuyentes -con modelización de la evasión y la elusión fiscal- como el funcionamiento de la administración tributaria. A continuación. en este tema. se exponen los contenidos fundamentales que confornrnn la TIO en la actualidad.

\subsubsection{Dimensión conductual: Incentivos} 
Las distorsiones impositivas, que son inevitables en cualquier sistema tributario, alteran los comportamientos que hubieran tenido lugar, bajo un principio de racionalidad económica, en  ausencia de impuestos, dando lugar a un coste de eficiencia cuya valoración monetaria identificamos con el concepto de exceso de gravamen. El exceso de gravamen representa la pérdida neta de bienestar económico causada por un impuesto distorsionante en comparación con un escenario hipotético en el que la recaudación obtenida fuese aportada por un impuesto de suma fija o capitación (lump sum), que no generase ninguna clase de efecto sustitución. Desde el punto de vista Cl:lantitativo. en una primera aproximación, el exceso de gravamen es la diferencia entre el valor monetario de los excedentes del consumidor y del productor en un mercado sin de impuestos y los valores de ambos tras la introducción de los impuestos, una vez descontada la recaudación. Cuanto mayor sea el exceso de gravamen, mayor será la pérdida de eficiencia económica que resulta de la imposición. A medida que nos adentremos en este tema, exploraremos cómo los principios de neutralidad, las distorsiones impositivas y el exceso de gravamen están interconectados y cómo afectan la toma de decisiones económicas, la asignación de recursos y el bienestar general de una sociedad. Comprender estos conceptos es esencial para disefiar políticas fiscales eficientes y equitativas, así como para evaluar criticamente el impacto de las reformas tributarias en la economía.


\subsection{Trilema de la Economía Pública: ¿Qué compromiso buscamos?}

Como hemos visto, la imposición óptima depende de la disyuntiva entre «eficiencia» y «equidad». Sin embargo, el uso de estos conceptos en la teoría de la imposición óptima no siempre se corresponde estrechamente con su uso común. En el contexto de la teoría de la imposición óptima, un impuesto equitativo es aquel que garantiza una distribución socialmente deseable de la carga fiscal; un impuesto eficiente es aquel con un pequeño exceso de gravamen. En el debate público, por otro lado, un impuesto equitativo es a menudo aquel que impone obligaciones iguales a las personas que tienen la misma capacidad de pago, y un sistema fiscal eficiente es aquel que mantiene bajos los gastos administrativos y de cumplimiento. Estas nociones alternativas de equidad y eficiencia en la imposición son el tema de esta sección.

Equidad horizontal

El humorista estadounidense Will Rogers dijo una vez: «La gente quiere impuestos justos más de lo que quiere impuestos más bajos. Quiere saber que cada hombre está pagando su parte proporcional de acuerdo con su riqueza». Este criterio para evaluar un sistema fiscal está incorporado en la noción de equidad horizontal de los economistas: las personas en posiciones iguales deberían ser tratadas por igual. Para hacer de la equidad horizontal una idea operativa, se debe definir qué son «posiciones iguales». Rogers sugiere la riqueza como un índice de capacidad de pago, pero también podrían utilizarse la renta y el gasto.

Desafortunadamente, todas estas medidas representan los resultados de las decisiones de las personas y no son realmente medidas adecuadas de una posición igual. Considérense dos individuos, ambos capaces de ganar 10 dólares por hora. El Sr. A elige trabajar 1.500 horas cada año, mientras que la Sra. B trabaja 2.200 horas cada año. La renta de A es de 15.000 dólares y la de B es de 22.000 dólares, de modo que, en términos de renta, A y B no están en «posiciones iguales». Sin embargo, en un sentido importante, A y B son iguales, porque sus capacidades de obtención de ingresos son idénticas; simplemente sucede que B trabaja más. Así, debido a que el esfuerzo laboral está, al menos hasta cierto punto, bajo el control de las personas, dos individuos con rentas diferentes pueden estar en realidad en posiciones iguales. Una crítica similar se aplicaría al gasto o a la riqueza como criterio para medir posiciones iguales.

Estos argumentos sugieren que se considere el salario por hora del individuo, en lugar de la renta, como candidato para medir posiciones iguales, pero esta idea también tiene problemas. En primer lugar, las inversiones en capital humano —educación, formación en el puesto de trabajo y atención sanitaria— pueden influir en el salario por hora. Si el Sr. A tuvo que ir a la universidad para ganar el mismo salario que la Sra. B es capaz de ganar únicamente con un título de educación secundaria, ¿es justo tratarlos de la misma manera? En segundo lugar, calcular el salario por hora requiere dividir los ingresos totales por las horas de trabajo, pero estas últimas no son fáciles de medir. (¿Cómo debería contabilizarse el tiempo dedicado a consultar Facebook?) De hecho, para una renta dada, a un trabajador le resultaría conveniente exagerar sus horas de trabajo para poder declarar un salario por hora más bajo y pagar menos impuestos. Presumiblemente, se podría inducir a los jefes a colaborar con sus empleados a cambio de una parte del ahorro fiscal.

Como alternativa a medir una posición igual ya sea en términos de rentas o de salarios por hora, Feldstein [1976] sugiere que se defina en términos de utilidades. De ahí, la definición de la equidad horizontal basada en la utilidad: (a) si dos individuos estarían igualmente bien (tendrían el mismo nivel de utilidad) en ausencia de imposición, también deberían estar igualmente bien si existe imposición; y (b) los impuestos no deberían alterar el ordenamiento de la utilidad: si A está mejor que B antes de impuestos, debería estar mejor después.

Para evaluar las implicaciones de la definición de Feldstein, supongamos primero que todos los individuos tienen las mismas preferencias, es decir, funciones de utilidad idénticas. En este caso, los individuos que consumen las mismas mercancías (incluido el ocio) deberían pagar el mismo impuesto o, de manera equivalente, todos los individuos deberían enfrentarse al mismo esquema impositivo. De lo contrario, los individuos con niveles de utilidad iguales antes de impuestos tendrían utilidades diferentes después de impuestos.

Supongamos ahora que las personas tienen gustos diversos. Por ejemplo, supongamos que hay dos tipos de individuos, Gourmets y Bañistas. Ambos grupos consumen alimentos (que se compran utilizando la renta) y ocio, pero los Gourmets otorgan un valor relativamente alto a los alimentos, al igual que los Bañistas al tiempo de ocio. Supongamos además que, antes de cualquier imposición, los Gourmets y los Bañistas tienen niveles de utilidad idénticos. Si se impone a todo el mundo el mismo impuesto proporcional sobre la renta, los Gourmets necesariamente quedan en peor situación que los Bañistas, porque los primeros necesitan cantidades relativamente grandes de renta para mantener sus hábitos alimentarios. Así, aunque este impuesto sobre la renta es perfectamente equitativo juzgado según la definición tradicional de equidad horizontal, no es equitativo según la definición basada en la utilidad. De hecho, mientras los gustos por el ocio difieran, cualquier impuesto sobre la renta viola la definición de equidad horizontal basada en la utilidad.

Por supuesto, las dificultades prácticas que implica medir las utilidades de los individuos excluyen la posibilidad de tener un impuesto sobre la utilidad. No obstante, la definición de equidad horizontal basada en la utilidad tiene algunas implicaciones provocadoras para la política. Supongamos de nuevo que todos los individuos tienen las mismas preferencias. Entonces puede demostrarse que cualquier estructura fiscal existente no viola la definición de equidad horizontal basada en la utilidad si los individuos son libres de elegir sus actividades y gastos.

Para ver por qué, supongamos que en un tipo de trabajo una gran parte de la remuneración consiste en ventajas que no están sujetas a impuestos —oficinas agradables, acceso a una piscina, etcétera—. En otra ocupación, la remuneración es exclusivamente monetaria, y toda ella está sujeta al impuesto sobre la renta. Según la definición tradicional, esta situación constituye una violación de la equidad horizontal, porque una persona en el trabajo con muchas ventajas tiene una carga fiscal demasiado pequeña. Pero, si ambos acuerdos coexisten y los individuos son libres de elegir, entonces las recompensas netas después de impuestos (incluidas las ventajas) deben ser las mismas en ambos trabajos. ¿Por qué? Supongamos que la recompensa neta después de impuestos es mayor en los trabajos con ventajas. Entonces los individuos migran hacia estos trabajos para aprovecharlas. Pero el aumento de la oferta de trabajadores en estos empleos reduce sus salarios. El proceso continúa hasta que los rendimientos netos son iguales. En resumen, aunque las personas en las distintas ocupaciones pagan impuestos desiguales, no existe inequidad horizontal debido a los ajustes en el salario antes de impuestos.

Algunos sugieren que ciertas ventajas fiscales disponibles únicamente para los ricos son fuentes de inequidad horizontal. Según la definición basada en la utilidad, esta noción es errónea. Si estas ventajas están abiertas a todo el mundo con una renta alta, y todas las personas de renta alta tienen gustos idénticos, entonces las ventajas pueden, en efecto, reducir la progresividad fiscal, pero no tienen efecto alguno sobre la equidad horizontal.

Llegamos a una conclusión sorprendente: dados unos gustos comunes, una estructura fiscal preexistente no puede implicar inequidad horizontal. Más bien, todas las inequidades horizontales surgen de cambios en las leyes fiscales. Esto se debe a que los individuos adquieren compromisos basados en las leyes fiscales existentes que son difíciles o imposibles de revertir. Por ejemplo, las personas pueden comprar casas más grandes debido al tratamiento fiscal preferente de las viviendas ocupadas por sus propietarios. Cuando se modifican las leyes fiscales, su bienestar disminuye y se viola la equidad horizontal. Como expresó un congresista: «Parece injusto para las personas que han hecho algo de buena fe cambiarles la ley».⁹ Estas observaciones dan un nuevo significado al dicho: «El único impuesto bueno es un impuesto antiguo».

El hecho de que los cambios fiscales puedan generar inequidades horizontales no implica necesariamente que no deban llevarse a cabo. Después de todo, los cambios fiscales pueden mejorar la eficiencia y/o la equidad vertical. Sin embargo, los argumentos sugieren que podría ser apropiado facilitar la transición al nuevo sistema fiscal. Por ejemplo, si se anuncia que una determinada reforma fiscal no entrará en vigor hasta unos años después de su aprobación, las personas que han basado su comportamiento en la antigua estructura fiscal podrán realizar al menos algunos ajustes al nuevo régimen. El problema de encontrar procesos justos para cambiar los regímenes fiscales —conocido como equidad transitoria— es muy difícil, y no hay muchos resultados disponibles sobre el tema.

Las implicaciones muy conservadoras de la definición de equidad horizontal basada en la utilidad no deberían causar gran sorpresa, porque implícita en la definición está la noción de que el statu quo antes de impuestos tiene una validez ética especial. (De lo contrario, ¿por qué preocuparse por los cambios en el ordenamiento de las utilidades?) Sin embargo, no resulta en absoluto obvio por qué el statu quo merece ser defendido. Una característica más general de la definición basada en la utilidad es su enfoque en los resultados de la imposición. En contraste, algunos han sugerido que la esencia de la equidad horizontal consiste en imponer restricciones a las reglas que rigen la selección de los impuestos, en lugar de proporcionar criterios para juzgar sus efectos. Así, la equidad horizontal excluye los impuestos caprichosos, o los impuestos basados en características irrelevantes. Por ejemplo, podemos imaginar que el gobierno imponga impuestos especiales de suma fija a las personas pelirrojas, o que aplique impuestos muy diferentes a los pasteles de ángel y a los pasteles de chocolate. La definición de equidad horizontal basada en reglas presumiblemente excluiría tales impuestos de consideración, incluso si tuvieran efectos deseables sobre la eficiencia o la distribución. En este sentido, las disposiciones de la Constitución de Estados Unidos que excluyen ciertos tipos de impuestos pueden interpretarse como un intento de garantizar la equidad horizontal. (Véase el Capítulo 1.)

Sin embargo, identificar el conjunto permisible de características en las que basar la imposición constituye un problema. La mayoría de las personas estarían de acuerdo en que la religión y la raza deberían ser irrelevantes a efectos de determinar la obligación tributaria. Por otro lado, existe un desacuerdo considerable acerca de si el estado civil debería o no influir en las cargas fiscales (véase el Capítulo 17). E incluso existiendo acuerdo en que ciertas características son bases legítimas para la discriminación, sigue existiendo el problema de cuánta discriminación es apropiada. Todo el mundo está de acuerdo en que una discapacidad física grave debería tenerse en cuenta al determinar la obligación tributaria personal. Pero ¿hasta qué punto debe ser mala su visión para tener derecho a un tratamiento fiscal especial como persona ciega? ¿Y en qué cantidad debería reducirse su factura fiscal?

Nos vemos obligados a concluir que la equidad horizontal, sea cual sea su definición, es un concepto bastante amorfo. Sin embargo, tiene un enorme atractivo como principio de diseño fiscal. Las nociones de justicia entre iguales, independientemente de su vaguedad, continuarán desempeñando un papel importante en el desarrollo de la política fiscal.


\textcolor{red}{\textbf{Origen}: GRUBER}
Con frecuencia se consideran dos objetivos distributivos al medir la equidad fiscal. El primero es la equidad vertical, el principio de que los grupos con más recursos (mayor renta, mayor riqueza, mayores beneficios, etc.) deberían pagar impuestos más altos que los grupos con menos recursos. Esta idea está relacionada con el concepto de equidad analizado en los Capítulos 2 y 17, que se refería a la distribución de recursos entre grupos de renta (o capacidad) más alta y más baja. Las preocupaciones por la equidad vertical en la imposición podrían estar motivadas, por ejemplo, por una función de bienestar social utilitarista que exige la redistribución desde los grupos de menor utilidad marginal del consumo hacia los de mayor utilidad marginal del consumo en la sociedad.

Otro concepto de equidad que se plantea a menudo en los debates sobre política fiscal es la equidad horizontal, el principio de que los individuos que son similares pero que toman diferentes decisiones económicas o de estilo de vida deberían ser tratados de la misma manera por el sistema fiscal. Considérense dos sistemas estatales de impuestos sobre las ventas. Un estado establece un impuesto sobre las ventas del 5 % sobre todos los bienes. Otro estado implementa un sistema mediante el cual, cada vez que usted realiza una compra, el cajero lanza una moneda, y usted no paga ningún impuesto sobre las ventas si sale cara y paga un impuesto sobre las ventas del 10 % si sale cruz. Este último sistema recaudará, en promedio, la misma cantidad de dinero que el primero, pero es mucho menos equitativo horizontalmente porque dos individuos idénticos podrían acabar pagando impuestos muy diferentes.

Este ejemplo extremo ilustra claramente una inequidad horizontal, pero en realidad las inequidades horizontales son difíciles de definir. Imaginemos que mi amigo y yo somos idénticos en términos de inteligencia, educación y motivación. Yo elijo pasar más tiempo en casa con mis hijos, mientras que mi amigo elige dedicar más tiempo a su trabajo. Aunque somos iguales en muchos aspectos, mi amigo tiene una renta más alta que yo y, como resultado, pagará impuestos sobre la renta más altos.

¿Viola este resultado la equidad horizontal? Por un lado, tenemos cantidades diferentes de renta, por lo que pagamos impuestos diferentes, lo cual parece horizontalmente equitativo. Por otro lado, somos dos personas idénticas en términos de capacidades y recursos subyacentes, pero, debido a las diferentes elecciones que hemos hecho, pagamos cantidades diferentes de impuestos, lo cual parece horizontalmente inequitativo.

Siempre que la cantidad de impuestos pagados dependa de las elecciones realizadas por los individuos, existirá tal dilema: individuos con recursos subyacentes idénticos que toman decisiones diferentes pagarán cantidades diferentes de impuestos. El único momento en que las inequidades horizontales son inequívocas es cuando los impuestos difieren por razones independientes de la elección, como en el ejemplo anterior de imposición aleatoria. Así, las violaciones de la equidad horizontal están, en última instancia, en el ojo de quien las contempla, un hecho desafortunado porque las preocupaciones por la equidad horizontal se plantean constantemente en los debates fiscales y a menudo se distorsionan para ajustarse a las opiniones de los defensores o de los opositores de una determinada propuesta fiscal.

Medición de la equidad vertical

Aunque la equidad horizontal es a menudo difícil de definir y medir, existen medidas más estándar de equidad vertical que son fundamentales en los debates sobre política fiscal. La mayoría de los analistas concluyen que, para ser verticalmente equitativos, los sistemas fiscales deben ser progresivos: los tipos impositivos medios efectivos deben aumentar con la renta, de modo que los ricos paguen en impuestos una proporción mayor de su renta que los pobres. (Por ejemplo, un sistema fiscal progresivo sería aquel en el que los individuos pagan el 10 % de su renta en impuestos con una renta de 10.000 dólares, pero pagan el 30 % de su renta en impuestos con una renta de 100.000 dólares). Los sistemas fiscales en los que el tipo impositivo medio efectivo no cambia con la renta son sistemas fiscales proporcionales porque todo el mundo paga la misma proporción de su renta en impuestos. (Por ejemplo, los individuos pagan el 15 % de la renta en impuestos independientemente de si ganan 10.000 o 100.000 dólares). Los sistemas fiscales en los que los tipos impositivos medios efectivos disminuyen con la renta son sistemas fiscales regresivos. (Por ejemplo, los individuos pagan el 15 % de su renta en impuestos con una renta de 10.000 dólares, pero pagan solo el 10 % de su renta en impuestos con una renta de 100.000 dólares).

El proceso político de medir la equidad fiscal¹⁴

Como sugiere la discusión anterior, medir la equidad fiscal puede ser difícil. Existen varias formas diferentes de medir la equidad, y es probable que los políticos elijan aquella que mejor se ajuste a sus agendas al defender u oponerse a un cambio fiscal. Un excelente ejemplo de este proceso es el debate sobre las reducciones del impuesto sobre la renta propuestas por el presidente George W. Bush y promulgadas por el Congreso en 2003. Estas reducciones fiscales aceleraron reducciones ya programadas de los tipos del impuesto sobre la renta, ampliaron las desgravaciones fiscales para las parejas casadas, aumentaron el crédito pagado a las familias con hijos y aumentaron las desgravaciones fiscales para las empresas.

Los críticos demócratas se opusieron a estas reducciones fiscales por motivos de «equidad». Señalaron, por ejemplo, que el 44 % de las reducciones fiscales de esta ley iría al 1 % superior de los contribuyentes. La administración Bush reconoció ese hecho, pero respondió señalando que estos contribuyentes de la parte superior ya pagan el 38 % de todos los impuestos sobre la renta. Por tanto, esta reducción de su factura fiscal era aproximadamente proporcional a sus pagos existentes del impuesto sobre la renta. Así, desde el punto de vista de los defensores de la ley, esta era una reducción justa de los impuestos para quienes actualmente pagan más en impuestos.

Los demócratas respondieron destacando que, aunque el 1 % superior de los contribuyentes paga el 38 % de los impuestos sobre la renta, paga solo el 30 % de todos los impuestos. Debido a que el tipo del impuesto sobre las nóminas es fijo en lugar de aumentar, y debido a que la base imponible de los impuestos de Seguridad para la Vejez e Ingresos por Discapacidad (OASDI) está limitada para quienes tienen ingresos elevados, nuestro sistema de impuestos sobre las nóminas es menos progresivo que nuestro sistema de impuesto sobre la renta. Por tanto, el 1 % superior estaba obteniendo una desgravación fiscal (el 44 % de la reducción fiscal) que estaba muy fuera de proporción con respecto a su participación en los pagos totales (el 30 % del total de impuestos pagados). Esto, sostenían los demócratas, era injusto.

La administración contraatacó señalando que 34 millones de familias con hijos recibirían una reducción fiscal media de 1.549 dólares cada una. Pero los críticos señalaron que este era un uso engañoso de la palabra «media». Estas cifras medias estaban infladas por el hecho de que la parte del león de la reducción fiscal corresponde a los hogares con las rentas más altas. Como lo expresó el economista y columnista del New York Times Paul Krugman: «Cuando Bill Gates entra en un bar, el patrimonio neto medio de los clientes se dispara, pero eso no convierte a todos los que están en el bar en multimillonarios».¹⁵ Aunque es cierto que 34 millones de familias con hijos obtendrían una reducción fiscal media de 1.549 dólares, esta media estaba compuesta tanto por 10 millones de familias que recibirían una reducción fiscal inferior a 100 dólares como por 200.000 familias (con rentas superiores a 1 millón de dólares al año) que recibirían una reducción fiscal de 93.500 dólares. Las familias situadas en el centro de la distribución de la renta recibieron una reducción fiscal media de solo 217 dólares.

Como se dice a menudo en Washington, D. C., la posición que uno adopta sobre una cuestión depende de dónde esté sentado. Para la mayoría de los republicanos, esta reducción fiscal recompensaba justamente a quienes soportaban la mayor carga del impuesto sobre la renta vigente. Para la mayoría de los demócratas, esta reducción fiscal recompensaba injustamente a los ricos de manera desproporcionada respecto de su carga fiscal total vigente. Ninguna de las partes del debate hizo realmente hincapié en la medida generalmente preferida por los economistas de los efectos distributivos de las políticas fiscales, que es cómo afectan a la distribución de la renta después de impuestos. Las reformas fiscales progresivas estrecharán la distribución de la renta después de impuestos; las reformas fiscales regresivas la ampliarán. Según las evaluaciones del Tax Policy Center, de carácter no partidista, tanto de las reducciones fiscales de 2001 como de las de 2003, el quintil inferior de la distribución de la renta vio aumentar sus rentas después de impuestos en un 0,7 % como resultado de estas reducciones fiscales, mientras que el quintil superior vio aumentar sus rentas después de impuestos en un 4,4 %, y el 0,1 % superior de los contribuyentes vio aumentar sus rentas después de impuestos en un 7,5 %. Claramente, según esta medida, los cambios fiscales fueron altamente regresivos, dando lugar a una ampliación de la distribución de las rentas después de impuestos.¹⁶








El sistema de ingresos de Estados Unidos está bajo ataque. Los críticos sostienen que es ineficiente, injusto e indebidamente complicado. Pero cuando estos críticos ofrecen propuestas de reforma, sus ideas son generalmente atacadas por las mismas razones. ¿Cómo debemos elegir? Nuestro objetivo en este capítulo es establecer un conjunto de criterios para evaluar los sistemas tributarios del mundo real. Comenzamos examinando consideraciones de eficiencia y distribución que encajan perfectamente dentro del marco de la economía del bienestar convencional. Luego pasamos a otros criterios que no encajan tan claramente, pero que, sin embargo, tienen una importancia y un atractivo considerables.


\clearpage
\section{Imposición Óptima: Marco analítico}
\puntosclave{
    Un sistema impositivo, con diversas figuras tributarias, debe minimizar el exceso de gravamen global. La cuantía, diseño y número de tipo dependerán, fundamentalmente, de: (i) el importe de recaudación necesaria; (ii) las elasticidades compensadas de la demanda y oferta de los bienes; y, (iii) los precios y cantidades intercambiadas bajo un marco de no intervención.

    Si se opta por obtener la recaudación buscada gravando solo un bien de consumo, la mejor opción en términos de neutralidad es gravar el bien con menor elasticidad de demanda compensada. En cambio, si se opta por gravar el consumo de más de un bien, la mejor opción es gravarlos todos, con tipos diferenciados que igualen el exceso de gravamen marginal generado por cada euro adicional de recaudación aportado por cada bien.
    
    Si se opta por gravar más de un bien y los efectos renta y sustitución cruzados son nulos, los tipos de gravamen óptimos deben ser inversamente proporcionales a las elasticidades precio de las respectivas demandas compensadas. En cambio, si los bienes no son independientes, los tipos óptimos deben reducir las demandas compensadas de todos los bienes gravados en la misma proporción.

    Si se opta por gravar el consumo de más de un bien, y el grado de complementariedad o sustituibilidad respecto del ocio es el mismo para todos bienes, el gravamen impositivo uniforme (con un único tipo) sobre el consumo es óptimo.

    La incorporación de la justicia distributiva como principio a tener en cuenta en la búsqueda de sistema impositivo óptimo sitúa su diseño ante el conflicto entre eficiencia y equidad. 

    La optimalidad de un sistema impositivo redistributivo solo es aceptable si el aumento del bienestar social que se produce por la reducción de la desigualdad de las utilidades individuales es mayor que la provocada por las distorsiones generadas por su aplicación.
    
    La progresividad óptima de un impuesto lineal sobre la renta exige fijar un tipo marginal óptimo mayor cuánto: (i) mayor sea la aversión a la desigualdad de la sociedad; (ii) mayor sea la recaudación deseada; y, (iii) menor sea la elasticidad-salario promedio de la oferta de trabajo compensada. Por el contrario, en un impuesto sobre la renta con varios tipos, el tipo marginal óptimo será mayor cuánto: (i) cuanto mayor sea la aversión a la desigualdad de la sociedad; (ii) mayor sea la recaudación deseada; y, (iii) menor sea la elasticidad-salario de la oferta de trabajo compensada del individuo $i$-ésimo.
    
    Cuando se pueden combinar la imposición progresiva sobre la renta personal como el gravamen del consumo de bienes con tipos diferenciados, si los individuos solamente se diferencian en su capacidad de obtención de renta (bajo el supuesto de que la utilidad individual es separable en la cantidad de trabajo ofertado y en el consumo de bienes), el diseño óptimo del impuesto directo hace innecesaria la utilización del impuesto indirecto con fines redistributivos.

    Con otros supuestos, razonables en un contexto real de aplicación, como son la existencia de evasión fiscal o la heterogeneidad de los individuos más allá de su capacidad para obtener rentas, el diseño de sistemas fiscales óptimos requiere contar simultáneamente con impuestos directos e indirectos. En cambio, no existen reglas concluyentes sobre cuál debe ser la combinación óptima de ambas formas de imposición.
}

Como decíamos en la introducción, el campo de estudio de la Economía de la Imposición nos ha proporcionado las herramientas necesarias para pasar de la cuestión positiva de cómo medir el exceso de gravamen a la cuestión normativa de cómo debería tenerse en cuenta la existencia de éste en el diseño del sistema fiscal, con un objetivo recaudatorio definido ($\bar T$). 

Abordamos esta cuestión normativa con referencia a dos tipos diferentes de impuestos: la imposición sobre los bienes y la imposición sobre la renta.

\subsection{Teoría de la Imposición Óptima: Imposición Indirecta}
La teoría de la imposición óptima sobre los bienes comenzó con el economista de principios del siglo XX Frank RAMSEY, quien consideró el problema de un gobierno con una necesidad presupuestaria dada y formuló el problema de la imposición óptima planteando la pregunta: ¿cómo podemos recaudarla con la menor cantidad de distorsión? En otras palabras, ¿cómo debería un gobierno establecer sus tipos impositivos para minimizar la pérdida irrecuperable de eficiencia del sistema fiscal al tiempo que satisface su necesidad presupuestaria?

\subsubsection{Base única: }





\subsubsection{Base general: }
Un impuesto sobre un bien de estas características puede parecer macabro, incluso absurdo. Sin embargo, la realidad es que, en términos de eficiencia, resulta prácticamente equivalente a un impuesto de suma fija. 

Todos los sistemas fiscales gravan un conjunto determinado de bienes de consumo, no solo uno. 

En Florida, las facturas de telefonía inalámbrica están gravadas a una tasa del $16,23\%$; la mayoría de los demás bienes (excepto los alimentos, que están exentos) están gravados a una tasa del $6\%$ por ciento. ¿Debería el servicio de telefonía inalámbrica estar gravado a una tasa más alta que otras cosas? Este es solo un ejemplo de una cuestión de política económica muy general y muy importante: ¿A qué tasas deberían gravarse los distintos bienes y servicios? El propósito de la teoría de la imposición óptima sobre los bienes es proporcionar un marco para responder a esta pregunta.\footnote{%
    Por supuesto, no podemos encontrar el conjunto correcto de impuestos sin conocer el objetivo del gobierno. Suponemos, de base, que el único objetivo es financiar los gastos del Estado con el menor exceso de gravamen y sin utilizar ningún impuesto de suma fija.
}

Consideremos la situación de Ana, una ciudadana que consume solo dos bienes, $X$ e $Y$, además de ocio, $l$. El precio de $X$ es $P_X$, el precio de $Y$ es $P_Y$, y la tasa salarial (precio del ocio) es $w$. El número máximo de horas por año que Ana puede trabajar está fijado en $\bar T$, la cantidad de tiempo que queda después de dormir. Se deduce que las horas de trabajo son $(\bar T − l)$: todo el tiempo que no se dedica al ocio se dedica al trabajo. El ingreso es el producto de la tasa salarial y las horas de trabajo: $w(\bar T − l)$. Suponiendo que Ana gasta todo su ingreso en los bienes $X$ e $Y$ (no hay ahorro), su restricción presupuestaria es:

\eqblock{
\begin{aligned}
    &\text{Sin impuestos}: \, \underset{{\scriptsize\text{Ingresos Totales}}}{w(\bar T-l)}=\underset{{\scriptsize\text{Gastos Totales}}}{P_X\,X+P_Y\,Y}\quad\overset{\scriptsize{\text{reescribimos}}}{\Longrightarrow}\quad \underset{{\scriptsize\substack{\text{Valor dotación}\\\text{de tiempo}}}}{w\,\bar T}=P_X\,X+P_Y\,Y+w\,l\\[2em]
    &\text{Con impuestos}: \, w\bar T = (1 + t)P_XX + (1 + t)P_YY + (1 + t)wl\\[0.5em]
    &\hspace{7em}\quad\overset{\scriptsize{\text{dividimos por $(1+t)$}}}{\Longrightarrow}\quad \frac{1}{1+t}\,w\bar T = P_XX + P_YY + wl
\end{aligned}
}{Restricción presupuestaria (sin tributación). Imposición óptima: Bienes}

Suponiendo que es posible gravar $X$, $Y$ y $l$ a la misma tasa ad valorem, $t$, vemos como el impuesto eleva el precio efectivo de $X$ a $(1 + t)P_X$, el de $Y$ a $(1 + t)P_Y$ y el de $l$ a $(1 + t)w$. Comparando entre las dos ecuaciones podemos señalar el siguiente hecho: un impuesto sobre todos los bienes, incluido el ocio, a la misma tasa porcentual, $t$, equivale a reducir el valor de la dotación de tiempo de $w\bar T \to [1/(1 + t)]w\bar T$. 

Sin embargo, dado que $w$ y $\bar T$ son fijos, su producto también es fijo; para cualquier valor de la tasa salarial, un individuo no puede cambiar el valor de su dotación de tiempo. Por lo tanto, un impuesto proporcional sobre la dotación de tiempo es, en efecto, un impuesto de suma fija y, como sabemos, no tienen exceso de gravamen. 

Por tanto, concluimos que un impuesto a la misma tasa sobre todos los bienes, incluido el ocio, equivale a un impuesto de suma fija y no tiene exceso de gravamen. Suena bien, ¿verdad? 

Hay un problema: imponer un impuesto sobre el tiempo de ocio es imposible. Los únicos instrumentos fiscales disponibles son los impuestos sobre los bienes $X$ e $Y$. Por lo tanto, cierto exceso de gravamen es generalmente inevitable. El objetivo de la imposición óptima sobre los bienes es seleccionar las tasas impositivas sobre $X$ e $Y$ de tal manera que el exceso de gravamen derivado de recaudar los ingresos fiscales requeridos sea lo más bajo posible. Podría parecer que la solución a este problema es gravar $X$ e $Y$ a la misma tasa, lo que se denomina imposición neutral. Pero, como veremos, la imposición neutral no es, en general, eficiente.

\paragraph{La regla de RAMSEY}
Entonces, para recaudar los ingresos con el menor exceso de gravamen posible, ¿cómo deberían establecerse las tasas impositivas sobre $X$ e $Y$? Para minimizar el exceso de gravamen total, el exceso de gravamen marginal del último euro de ingresos recaudado de cada bien debe ser el mismo. De lo contrario, sería posible reducir el exceso de gravamen total aumentando la tasa sobre el bien con el menor exceso de gravamen marginal y reduciendo al mismo tiempo la tasa sobre el bien con el mayor exceso de gravamen marginal.

Para explorar las consecuencias de este ejemplo típico de análisis marginal, supongamos, para simplificar, que para nuestra consumidora, $X$ e $Y$ son bienes independientes: no son ni sustitutos ni complementarios entre sí. Por lo tanto, un cambio en el precio de cualquiera de los bienes afecta a su propia demanda y no a la demanda del otro bien. Supongamos que Ana puede comprar todo el $X$ que quiera al precio $P_0$ (curva de oferta perfectamente elástica) y mostremos la demanda compensada de Ana de $X$, $D_X$:

\imagenfit{Fig1.png}{}{}

Supongamos que se aplica un impuesto unitario de $u_x$, lo que reduce la cantidad demandada de $X_0\to X_1:\Delta X$. Como sabemos, el exceso de gravamen del impuesto es el área del triángulo $abc$. Ahora supongamos que aumentamos el impuesto en $1$ unidad monetaria, de modo que pasa a ser $(u_1 + 1)$. El precio total es $P_0 + (u_x + 1)$; la cantidad demandada disminuye en $\Delta x$ hasta $X_2$; y el exceso de gravamen asociado es el triángulo $fec$. El exceso de gravamen marginal es la diferencia entre los dos triángulos, el trapecio $fbae$. El área del trapecio es la mitad de su altura ($\Delta_x$) multiplicada por la suma de sus bases $[u_x + (u_x + 1)]$. Así, el exceso de gravamen marginal es $\frac{1}{2}\Delta x[u_x + (u_x + 1)]$. Con un poco de álgebra, podemos simplificar esta expresión para obtener que el exceso de gravamen marginal es aproximadamente $\Delta X$:

\eqblock{
\begin{aligned}
    &\frac{1}{2}\Delta x[u_x+(u_x+1)]=\frac{1}{2}\Delta x(2u_x+1)=\Delta xu_x+\underset{{\scriptsize\text{pequeño, $\approx0$}}}{\frac{1}{2}\Delta x}\approx\Delta xu_x\\[1em]
    &\underset{{\scriptsize\text{los dos miden la pendiente de $D_x$}}}{\frac{1}{\Delta x}=\frac{u_x}{\Delta x}}\quad\Longrightarrow\quad \Delta xu_x=\boxed{\Delta X:\text{Exceso de gravamen marginal}}
\end{aligned}
}{Exceso de gravamen marginal: Derivación analítica (apoyo gráfico). Imposición óptima: Bienes}

Recordemos que la minimización del exceso de gravamen requiere información sobre el exceso de gravamen marginal del último euro de ingresos recaudado. Ahora que conocemos el exceso de gravamen marginal inducido por el aumento del impuesto, debemos calcular el aumento asociado de los ingresos. Entonces, todo lo que tenemos que hacer es dividir el exceso de gravamen marginal por el cambio en los ingresos. Por definición, este cociente es el exceso de gravamen marginal por euro incremental de ingresos recaudados.

Para calcular el cambio en los ingresos fiscales asociado con el aumento de la tasa de $u_x$ a $(u_x + 1)$, observemos que cuando la tasa impositiva es $u_x$, los ingresos fiscales son $u_xX_1$ (el impuesto por unidad multiplicado por el número de unidades vendidas). En la la representación gráfica anterior, esto es el rectángulo $hbaj$. De manera similar, cuando la tasa impositiva es $(u_x+ 1)$, los ingresos fiscales son $gfej$. Comparando estos dos rectángulos, vemos que cuando el impuesto aumenta, el gobierno gana el área $gfih$ pero pierde $ibae$. Así, el cambio en los ingresos es $gfih − ibae$. Utilizando álgebra:

\eqblock{
\begin{aligned}
    &(g,f,i,h) − (i,b,a,e)=X_2 − (X_1 − X_2)u_x\quad\Longrightarrow\quad \boxed{X_1 − \Delta X : \text{Ingreso fiscal marginal}}\\[2em]
    &X_1 − \Delta X\quad\sim\quad X_2(u_x+1)-X_1u_x=X_2+u_x(X_2-X_1)\\[0.5em]
    &\overset{{\scriptsize\text{sustituyendo $X_2=X_1-\Delta x$}}}{\Longrightarrow}\quad X_1-\Delta x-u_x\Delta x,\quad\text{pero}\quad\Delta x=\frac{\Delta X}{u_x}\\[0.5em]
    &\quad\Longrightarrow X_1-\frac{\Delta X(1+u_x)}{u_x}\quad\overset{{\scriptsize\text{como $u_x\gg 1$}}}{\Longrightarrow}\quad X_1-\Delta X
\end{aligned}
}{Ingreso fiscale marginal: Derivación analítica. Imposición óptima: Bienes}

Por tanto, poniendo todo en conjunto, obtenemos que el exceso de gravamen marginal por euro adicional de ingresos fiscales es:

\eqblock{
\begin{aligned}
    &\text{Regla de óptimo impositivo ($X$)}:\quad\boxed{\frac{\Delta X}{(X_1-\Delta X)}}\\[1em]
    &\text{Regla de óptimo impositivo ($X$)}:\quad\boxed{\frac{\Delta Y}{(Y_1-\Delta Y)}}\\[1em]
    &\text{Minimización}:\quad \frac{\Delta X}{(X_1-\Delta X)}=\frac{\Delta Y}{(Y_1-\Delta Y)}\quad\Longrightarrow\quad \boxed{\frac{\Delta X}{X_1}=\frac{\Delta Y}{Y_1}}
\end{aligned}
}{}

Exactamente el mismo razonamiento indica que si se aplica un impuesto unitario de $u_y$ sobre el otro bien, el exceso de gravamen marginal por último euro de ingresos al que se llega es el mismo; y, dado que la condición para minimizar el exceso de gravamen total es que el exceso de gravamen marginal por último euro de ingresos sea el mismo para cada bien, obtenemos la regla de fijación igualando las dos reglas de imposición óptima. Observemos que el cambio en una variable dividido por su valor total es simplemente el cambio porcentual de la variable. Por lo tanto, la última ecuación dice que, para minimizar el exceso de gravamen total, las tasas impositivas deberían establecerse de manera que la reducción porcentual de la cantidad demandada de cada bien sea la misma. 

Este resultado, denominado Regla de RAMSEY (1927), también se mantiene incluso en los casos en que $X$, $Y$ y $l$ son bienes relacionados, sustitutos o complementarios. Pero ¿por qué debería una imposición eficiente inducir cambios proporcionales iguales en las cantidades demandadas en lugar de cambios proporcionales iguales en los precios? Porque \textbf{el exceso de gravamen es una consecuencia de las distorsiones en las cantidades}. Minimizar el exceso de gravamen total requiere que todos estos cambios estén en la misma proporción.

El apéndice de este capítulo presenta las matemáticas de la elegante solución de Ramsey. Aquí analizamos la lección clave de su modelo: el gobierno debería establecer los impuestos entre las mercancías de modo que la razón entre la pérdida irrecuperable de eficiencia marginal y los ingresos marginales recaudados sea igual entre las mercancías:

Regla de Ramsey: establecer impuestos sobre las mercancías de tal manera que

MDWLᵢ / MRᵢ = λ

MDWL es la pérdida irrecuperable de eficiencia marginal derivada de aumentar el impuesto sobre el bien i, MR son los ingresos marginales recaudados por ese aumento del impuesto, y λ es el valor de los ingresos públicos adicionales. Esta constante mide el valor de tener otro dólar en manos del gobierno en relación con su siguiente mejor uso en el sector privado. Si λ es grande, implica que los ingresos públicos adicionales son bastante valiosos en relación con mantener el dinero en manos privadas; si λ es pequeño, entonces los ingresos públicos adicionales tienen poco valor en relación con el valor que los individuos privados atribuyen a tener ese dinero.⁷

Esta regla establece que la pérdida irrecuperable de eficiencia por dólar de ingresos fiscales asociada con un dólar adicional de impuestos sobre la mercancía i debería ser igual para todas las mercancías. Si el impuesto sobre el bien A tiene una MDWL/MR que es mayor que la MDWL/MR derivada de gravar el bien B, gravar el bien A causa más ineficiencia por dólar de ingresos recaudados que gravar el bien B. Recuérdese que MDWL es una función positiva del tipo impositivo; como se analizó anteriormente, unos impuestos más altos conducen a una mayor pérdida irrecuperable de eficiencia marginal porque alejan el mercado más del equilibrio competitivo (la pérdida irrecuperable de eficiencia está determinada por el cuadrado del tipo impositivo). Por lo tanto, para minimizar la ineficiencia en el mercado, el gobierno debería reducir la imposición sobre el bien A, reduciendo así su MDWL, y aumentar el impuesto sobre el bien B, incrementando su MDWL. Estos ajustes deberían continuar hasta que las razones MDWL/MR para ambos bienes sean iguales a λ, de modo que ambos bienes tengan el mismo coste de eficiencia por dólar de ingresos recaudados.

Si λ es grande, entonces los recursos adicionales para el gobierno tienen un valor elevado, por lo que la MDWL/MR debería ser grande para todos los impuestos sobre las mercancías (los tipos impositivos deberían ser altos); si λ es pequeño, entonces los recursos adicionales para el gobierno tienen un valor bajo, por lo que la MDWL/MR debería ser pequeña para todos los impuestos sobre las mercancías (los tipos impositivos deberían ser bajos). En otras palabras, el gobierno debería estar dispuesto a tener impuestos potencialmente ineficientes (MDWL alta) cuando tiene grandes necesidades presupuestarias. Esto indica al gobierno que debe establecer el coste marginal de la imposición (MDWL/MR) igual a su beneficio marginal (λ).

En principio, la Regla de Ramsey puede indicarnos el nivel óptimo de imposición entre las mercancías. En la práctica, los analistas de políticas normalmente no disponen de un valor medido para λ, el valor de los ingresos adicionales para el gobierno, por lo que la Regla de Ramsey se utiliza normalmente cuando se habla de reforma fiscal y de los costes y beneficios de pasar de un patrón existente de impuestos sobre las mercancías a otro patrón que recauda la misma cantidad de ingresos. La aplicación de las páginas 336–339 analiza el uso de la Regla de Ramsey para orientar la reforma fiscal.

\paragraph{Una reinterpretación de la regla de RAMSEY: Regla de la elasticidad inversa}

Es útil explorar la relación entre la Regla de RAMSEY y las elasticidades de la demanda. Sea $\eta_x$ la elasticidad compensada de la demanda de $X$. Sea $t_x$ la tasa impositiva sobre $X$, esta vez expresada como una tasa ad valorem en lugar de un impuesto unitario. Ahora bien, por definición de un impuesto ad valorem, $t_x$ es el aumento porcentual del precio inducido por el impuesto. Por lo tanto, $t_x\eta_x$ es el cambio porcentual en el precio multiplicado por el cambio porcentual en la cantidad demandada cuando el precio aumenta un $1\%$. Esto es simplemente la reducción porcentual de la demanda de $X$ inducida por el impuesto. Definiendo $t_y$ y $\eta_y$ de manera análoga, $t_y\eta_y$ es la reducción proporcional de $Y$. 

La Regla de RAMSEY dice que, para minimizar el exceso de gravamen, estas reducciones porcentuales de la cantidad demandada deben ser iguales:

\eqblock{
\begin{aligned}
    \text{Regla de RAMSEY}:\quad t_x\eta_x=t_y\eta_y\\[1em]
    \overset{{\scriptsize\text{dividiendo por $t_yn\eta_x$}}}{\Longrightarrow}\quad \boxed{\text{Regla de la elasticidad inversa}:\frac{t_x}{t_y}=\frac{\eta_x}{\eta_y}}
\end{aligned}
}{Regla de la elasticidad inversa: Derivación analítica desde Regla de RAMSEY. Imposición óptima: Bienes}

La regla de la elasticidad inversa nos dice algo muy sencillo: siempre que los bienes no estén relacionados en el consumo, las tasas impositivas deberían ser inversamente proporcionales a las elasticidades. Es decir, cuanto mayor sea $\eta_y$ en relación con $\eta_x$, menor debería ser $t_y$ en relación con $t_x$. La eficiencia no requiere que todas las tasas se establezcan uniformemente.

La intuición detrás de la regla de la elasticidad inversa es sencilla. Los impuestos eficientes distorsionan las decisiones lo menos posible. Como ya vimos, el potencial de distorsión es mayor cuanto más elástica sea la demanda de un bien. Por lo tanto, una imposición eficiente requiere que se apliquen tasas impositivas relativamente altas a los bienes relativamente inelásticos.

\begin{specialblock}{\textbf{EXTRA: Tarifas óptimas de usuario}}
    Hasta ahora hemos supuesto que toda la producción tiene lugar en el sector privado. El único problema del gobierno es establecer las tasas impositivas que determinan los precios al consumidor. A veces, el propio gobierno es el productor de un bien o servicio. En tales casos, el gobierno debe elegir directamente una tarifa de usuario. Como de costumbre, nos gustaría determinar la mejor tarifa de usuario posible.
    
    Analíticamente, los problemas del impuesto óptimo y de la tarifa de usuario están estrechamente relacionados. En ambos casos, el gobierno establece el precio final pagado por los consumidores. En el problema del impuesto óptimo, esto se hace indirectamente mediante la elección de la tasa impositiva, mientras que en el problema de la tarifa óptima de usuario, se hace directamente.
    
    ¿Cuándo debería el gobierno elegir producir un bien en lugar de comprarlo al sector privado? La producción gubernamental puede ser apropiada cuando el uso de algún bien o servicio está sujeto a costes medios continuamente decrecientes: cuanto mayor sea el nivel de producción, menor será el coste por unidad. En tales circunstancias, es poco probable que el mercado del servicio sea competitivo. Una sola empresa puede aprovechar las economías de escala y suministrar toda la producción de la industria, al menos para una región de tamaño considerable. Este fenómeno se denomina a menudo monopolio natural. Algunos ejemplos son los puentes, la electricidad y la televisión por cable. En algunos casos, estos bienes son producidos por el sector privado y regulados por el gobierno (electricidad); en otros, son producidos por el sector público (puentes). Aunque aquí estudiamos la producción pública, muchas de las ideas importantes se aplican a la regulación de los monopolios privados.
    
    La Figura 16.2 mide la producción del monopolio natural, $Z$, en el eje horizontal, y los euros en el vertical. La curva de coste medio se denota $AC_Z$. Por supuesto, disminuye continuamente en todos los niveles relevantes de producción. Debido a que el coste medio es decreciente, el coste marginal debe ser menor que el coste medio. Por lo tanto, la curva de coste marginal ($MC_Z$), que muestra el coste incremental de proporcionar cada unidad de $Z$, se encuentra por debajo de $AC_Z$. La curva de demanda de Z está representada por $D_Z$. La curva de ingreso marginal asociada es $MR_Z$. Esta muestra el ingreso incremental asociado con cada nivel de producción de $Z$.
    
    Para ilustrar por qué los costes medios decrecientes conducen a menudo a la producción por parte del sector público o a la producción regulada por parte del sector privado, consideremos qué ocurriría si $Z$ fuera producido por un monopolista no regulado. Un monopolista que busca maximizar los beneficios produce hasta el punto en que el ingreso marginal es igual al coste marginal, el nivel de producción $Z_m$ en la Figura 16.3. El precio asociado, $P_m$, se obtiene subiendo hasta la curva de demanda, $D_Z$. Los beneficios del monopolio son iguales al producto del número de unidades vendidas por el beneficio por unidad y están representados geométricamente por el rectángulo de color claro.
    
    ¿Es eficiente el nivel de producción $Z_m$? Según la teoría de la economía del bienestar, la eficiencia requiere que el precio sea igual al coste marginal: el valor que las personas atribuyen al bien debe ser igual al coste incremental para la sociedad de producirlo. En $Z_m$, el precio es mayor que el coste marginal. Por lo tanto, $Z_m$ es ineficiente. Esta ineficiencia, junto con el hecho de que la sociedad puede no aprobar la existencia de los beneficios del monopolio, proporciona una posible justificación para que el gobierno se haga cargo de la producción de $Z$.
    
    La prescripción política obvia parece ser que el gobierno produzca hasta el punto en que el precio sea igual al coste marginal. En la Figura 16.3, la producción en la que $P = MC$ se denota $Z^*$, y el precio asociado es $P^*$. Sin embargo, existe un problema: en el nivel de producción $Z^*$, el precio es menor que el coste medio. El precio $P^*$ es tan bajo que la operación no puede cubrir sus costes y sufre pérdidas. La pérdida total es igual al producto del número de unidades vendidas, $Z^*$, por la pérdida por unidad, medida como la distancia vertical entre la curva de demanda y $AC_Z$ en $Z^*$. Geométricamente, la pérdida es el rectángulo de color más oscuro de la Figura 16.3.
    
    ¿Cómo debería afrontar el gobierno este dilema? Se han propuesto varias soluciones.
    
    Fijación de precios según el coste medio
    
    Por definición, cuando el precio es igual al coste medio, no hay ni beneficios ni pérdidas: la empresa simplemente alcanza el punto de equilibrio. La operación ya no tiene que preocuparse por un déficit. Geométricamente, esto corresponde a la intersección de las curvas de demanda y coste medio en la Figura 16.3, donde la producción es $Z_A$ y el precio es $P_A$. Sin embargo, obsérvese que $Z_A$ es menor que $Z^*$. Aunque la fijación de precios según el coste medio conduce a una mayor producción que en el nivel de maximización de beneficios, sigue estando por debajo de la cantidad eficiente.
    
    Fijación de precios según el coste marginal con impuestos de suma fija
    
    Cobrar $P = MC$ y compensar el déficit mediante la imposición de impuestos de suma fija. Cobrar $P = MC$ garantiza la eficiencia en el mercado de $Z$; financiar el déficit con impuestos de suma fija sobre el resto de la sociedad garantiza que no se generen nuevas ineficiencias al cubrir el déficit. Sin embargo, existen dos problemas con esta solución:
    
    En primer lugar, como se señaló anteriormente, los impuestos de suma fija generalmente no están disponibles. El déficit tiene que financiarse mediante impuestos distorsionadores, como los impuestos sobre la renta o sobre los bienes. Si es así, la distorsión debida al impuesto puede superar con creces la ganancia de eficiencia en el mercado de $Z$.
    
    En segundo lugar, existe una creencia generalizada de que la equidad exige que los consumidores de un servicio proporcionado públicamente paguen por él: el llamado principio de los beneficios recibidos. Si este principio se toma en serio, es injusto compensar el déficit mediante impuestos generales. Si la guardia costera me rescata de un mar tormentoso, ¿por qué deberías pagarlo tú?
    
    \textbf{Una solución de RAMSEY.} Hasta ahora hemos estado examinando una empresa gubernamental de forma aislada. Supongamos que el gobierno gestiona varias empresas y que, como grupo, no pueden perder dinero, aunque cualquier empresa individual sí puede hacerlo. Supongamos además que el gobierno quiere que la financiación proceda de los usuarios de los servicios producidos por las empresas. ¿En cuánto debería exceder la tarifa de usuario de cada servicio su coste marginal?
    
    ¿Te resulta familiar esta pregunta? Debería, porque es esencialmente la misma que el problema del impuesto óptimo. En efecto, la diferencia entre el coste marginal y la tarifa de usuario es simplemente el impuesto que el gobierno aplica al bien. Y, al igual que en el problema del impuesto óptimo, el gobierno tiene que recaudar una determinada cantidad de ingresos; en este caso, suficiente para que el grupo de empresas alcance el punto de equilibrio. La regla de RAMSEY proporciona la respuesta: establecer las tarifas de usuario de manera que las demandas de cada bien se reduzcan proporcionalmente. Este análisis, por cierto, ilustra una de las características positivas de la teoría económica. A menudo, un marco desarrollado para estudiar un problema puede aplicarse de manera fructífera a otro problema que parece ser bastante diferente.
    
    \textbf{Visión general}. De las diversas posibilidades para hacer frente a los monopolios naturales, ¿cuál ha elegido Estados Unidos? En la mayoría de los casos, tanto las empresas de propiedad pública como las empresas privadas reguladas han seleccionado la fijación de precios según el coste medio. Aunque la fijación de precios según el coste medio es ineficiente, probablemente sea un compromiso razonable. Tiene la virtud de ser bastante sencilla y se ajusta al popular principio de los beneficios recibidos. Sin embargo, algunos economistas sostienen que sería deseable una mayor dependencia de la fijación de precios de RAMSEY.
\end{specialblock}

\paragraph{Regla de CORLETT-HAGUE}
CORLETT y HAGUE (1953) demostraron una implicación interesante de la regla de RAMSEY: cuando hay dos bienes, una imposición eficiente requiere gravar a una tasa relativamente alta el bien que es complementario del ocio. Para comprender intuitivamente este resultado, recordemos que, si fuera posible gravar el ocio, podría obtenerse un resultado de primer óptimo: podrían recaudarse ingresos sin exceso de gravamen. 

Aunque las autoridades fiscales no pueden gravar el ocio, pueden gravar los bienes que tienden a consumirse conjuntamente con el ocio, reduciendo indirectamente la demanda de ocio. Si los videojuegos se gravan a una tasa muy alta, las personas compran menos de ellos y dedican menos tiempo al ocio. En efecto, entonces, los impuestos elevados sobre los complementos del ocio proporcionan una manera indirecta de alcanzar el ocio y, por lo tanto, acercarse al resultado perfectamente eficiente que sería posible si el ocio pudiera gravarse.





\subsection{Teoría de la Imposición Óptima: Imposición Indirecta}
Esta formulación de elasticidad inversa del modelo de Ramsey pone de relieve las implicaciones bastante desagradables para la equidad del enfoque de Ramsey. Imaginemos que el gobierno tuviera solo dos bienes que pudiera gravar, cereales y caviar. La elasticidad de la demanda de caviar es mucho mayor que la de los cereales, por lo que la regla de la elasticidad inversa sugeriría que el gobierno gravara los cereales mucho más que el caviar. Esto significaría imponer a un bien consumido exclusivamente por grupos de renta más alta un impuesto mucho menor que el impuesto impuesto a un bien consumido por todos. Este resultado, aunque eficiente, podría violar el sentido de equidad fiscal de un gobierno entre grupos de renta (equidad vertical).

Un marco de imposición óptima sobre las mercancías puede abordar las preocupaciones de equidad teniendo en cuenta no solo la elasticidad de cada mercancía, sino también la distribución de la renta de sus consumidores. Los bienes que son consumidos desproporcionadamente por consumidores de renta más alta podrían tener un tipo impositivo superior al que implica la regla de la elasticidad inversa, y los bienes que son consumidos desproporcionadamente por consumidores de renta más baja podrían tener un tipo impositivo inferior al que implica la regla de la elasticidad inversa. Cuánto debería hacerse de esta «reponderación» de los impuestos óptimos entre las mercancías es una función del grado en que los gobiernos quieran intercambiar eficiencia por equidad. A medida que el gobierno se aleja de la regla de eficiencia de Ramsey incorporando cuestiones de equidad, el sistema fiscal se vuelve menos eficiente pero más equitativo.

Quizá debido a estas preocupaciones distributivas, en Estados Unidos se recurre relativamente poco a la imposición sobre las mercancías. La mayor parte de nuestros ingresos fiscales procede de gravar las rentas individuales; por tanto, analicemos ahora la cuestión de cómo diseñar sistemas en los que las obligaciones tributarias se basen en las rentas de las personas. 

Para enmarcar adecuadamente las cuestiones que trataremos, consideremos el debate acerca de la progresividad del impuesto de la renta. Los partidarios de la idea defienden que es una condición imprescindible para alcanzar la equidad; los opositores argumentan que es injusta e ineficiente. ¿Hasta qué punto debería ser progresivo el impuesto sobre la renta?\footnote{%
    El economista del siglo XIX, McCULLOCH, que se oponía a la imposición progresiva, sostuvo que una vez que se abandona la imposición proporcional, \textit{se está en el mar sin timón ni brújula, y no hay cantidad de injusticia y necedad que no se pueda cometer}. Como hemos visto, el objetivo de la teoría de la imposición óptima sobre la renta es proporcionar un timón, es decir, proporcionar una manera sistemática de pensar sobre el equilibrio correcto entre equidad y eficiencia.
}

\begin{specialblock}\textbf{Antecedentes: EDGEWORTH (1897)}
    A finales del siglo XIX, Edgeworth [1959/1897] examinó la cuestión de la imposición óptima sobre la renta utilizando un modelo sencillo basado en los siguientes supuestos: (recuperar del Tema 3.B.10)
    
    
    [...] Los supuestos de Edgeworth son prácticamente idénticos a los supuestos que subyacen al modelo de distribución óptima de la renta presentado en el Capítulo 12 bajo «Justificaciones de la redistribución de la renta». Allí mostramos que, con estos supuestos, la maximización del bienestar social requiere que la utilidad marginal de la renta de cada persona sea la misma. Cuando las funciones de utilidad son idénticas, las utilidades marginales son iguales solo si las rentas son iguales. Las implicaciones para la política tributaria son claras: los impuestos deberían establecerse de manera que la distribución de la renta después de impuestos sea lo más igualitaria posible. En particular, la renta debería tomarse primero de los ricos porque la utilidad marginal perdida es menor que la de los pobres. Si el gobierno requiere más ingresos incluso después de alcanzar la igualdad completa, la carga tributaria adicional debería distribuirse uniformemente.
    
    El modelo de Edgeworth, entonces, implica una estructura tributaria radicalmente progresiva: las rentas se nivelan desde arriba hasta alcanzar la igualdad completa. En efecto, los tipos impositivos marginales sobre los individuos de altos ingresos son del 100 por ciento. Sin embargo, como se destacó en el Capítulo 12, cada uno de los supuestos que subyacen a este análisis es cuestionable. En las últimas décadas, los economistas han investigado cómo cambian los resultados de Edgeworth cuando se relajan ciertos supuestos.
    
    Uno de los problemas más desconcertantes del análisis de Edgeworth es el supuesto de que la cantidad total de renta disponible para la sociedad es fija. Según este supuesto, los tipos impositivos confiscatorios no tienen ningún efecto sobre la cantidad de producción generada. De manera más realista, supongamos que las utilidades de los individuos dependen no solo de la renta, sino también del ocio. Entonces los impuestos sobre la renta distorsionan las decisiones laborales y crean excesos de gravamen (véase el Capítulo 15). Una sociedad con una función aditiva de bienestar social se enfrenta así a un dilema ineludible. Por una parte, desea asignar la carga tributaria para igualar la distribución de la renta después de impuestos. Sin embargo, en el proceso de hacerlo, reduce la cantidad total de renta real disponible. Un sistema óptimo de imposición sobre la renta —uno que maximiza el bienestar social— debe tener en cuenta los costes (en exceso de gravamen) de lograr una mayor igualdad. En el modelo de Edgeworth, el coste de obtener una mayor igualdad es cero, lo que explica la prescripción de un resultado perfectamente igualitario.
    
    ¿Cómo cambia el resultado de EDGEWORTH cuando se tienen en cuenta los incentivos al trabajo? 
\end{specialblock}

\subsubsection{Impuesto sobre la Renta Personal: Desarrollo analítico (MIRRLESS, 1971)}
\textcolor{red}{\textbf{OJO -> Aquí se debe reescribir el documento por la siguiente razón: el trabajo teórico más importante relativo a la tributación óptima sobre la renta es el de MIRRLESS (1971). Lo que MIRRLESS descubre es que, dada una FBS razonable, el esquema óptima es muy parecido a uno lineal. A partir de aquí, diversos estudios analizan la idoneidad de diferentes esquemas de tributación. STERN, por ejemplo, restringe el problema a los esquemas lineales y se pregunta ¿cuál sería el tipo óptimo?. Por su parte, relativo a los esquemas de tributación progresiva, son GRUBER y SAEZ (2002) los que aportan las conclusiones más interesantes que se presentan.}} 

El modelo de MIRRLEES (1971) constituye la contribución germinal de esta vertiente de la Teoría de la Imposición Óptima.

MIRRLEES (1971) analiza cómo los individuos toman decisiones sobre trabajo y ocio bajo la aplicación de impuestos progresivos, donde la productividad del trabajo desempeña un papel esencial a la hora de explicar las distorsiones provocadas por el gravamen de las rentas salariales obtenidas. El punto de partida del cual parte su desarrollo es sencillo: cuando la sociedad redistribuye recursos, probablemente reduce el tamaño total del pastel económico al reducir la generación de renta, mientras que al mismo tiempo iguala la distribución de las porciones del pastel.\footnote{%
    [EDGEWORTH] En el ejemplo anterior, no había ninguna disyuntiva; es decir, las rentas totales de la sociedad eran fijas, por lo que no había ningún coste de eficiencia derivado de aumentar los impuestos. Así, el sistema óptimo se centraba únicamente en la equidad, asegurando que todos los individuos tuvieran rentas iguales.
}
    
Dicho de otra forma, la tributación sobre la renta afecta al tamaño del pastel pues la tasa a la que se gravan las rentas generalmente determinará el tamaño de las rentas que están sujetas a tributación. Así, al diseñar impuestos óptimos sobre la renta, el gobierno necesita identificar un esquema de tasas impositivas entre los grupos de renta que maximice el bienestar social, reconociendo al mismo tiempo que aumentar las tasas impositivas tiene efectos contrapuestos (y, si los impuestos son suficientemente altos, negativos) sobre los ingresos. Es decir, de manera muy similar a la imposición óptima sobre bienes, representa un compromiso entre:\footnote{%
    Consideremos el ejemplo de un impuesto sobre la renta del trabajo. Los ingresos recaudados por este impuesto son iguales a la tasa impositiva multiplicada por la base imponible de los ingresos laborales. Supongamos que los trabajadores reducen la cantidad de trabajo que ofrecen al mercado a medida que cae su salario después de impuestos. Un aumento de la tasa impositiva sobre la renta del trabajo tendrá, por tanto, dos efectos sobre los ingresos fiscales. Primero, los ingresos fiscales aumentarán para un nivel dado de renta del trabajo. Segundo, sin embargo, en algún momento los trabajadores reducirán la cantidad de renta del trabajo que obtienen, y la base imponible se reducirá. Para tasas impositivas bajas, el primer efecto dominará: si partimos de una tasa impositiva de cero, no hay ninguna base imponible que pueda reducirse al aumentar la tasa impositiva, por lo que solo puede operar el primer efecto. Pero a medida que las tasas impositivas aumentan, el segundo efecto se volverá cada vez más importante: con una tasa impositiva del $100\%$, nadie trabajaría, la base sería cero y los impuestos no recaudarían ningún ingreso.
    
Estos dos efectos son la génesis de la famosa curva de Laffer, que fue un fundamento intelectual de las grandes reducciones de impuestos de principios de la década de 1980 en Estados Unidos. [\authorfont{Ver Tema 3.B.10}]
}
\begin{la}
    \item \textbf{Equidad vertical}. El bienestar social se maximiza cuando aquellos que tienen un nivel alto de consumo, y por tanto una utilidad marginal baja, son gravados más fuertemente, y aquellos que tienen un nivel bajo de consumo, y por tanto una utilidad marginal alta, son gravados menos fuertemente; aquellos con un consumo alto «echarán menos de menos el dinero» cuando este sea retirado mediante impuestos;
    \item \textbf{Respuestas conductuales}. A medida que los impuestos aumentan sobre cualquier grupo, los individuos de ese grupo pueden responder obteniendo menos renta. Esto significa que un aumento adicional de los impuestos recaudará menos ingresos porque la base de tributación es menor.
\end{la}


Al formular su modelo, MIRRLESS considera un escenario de información asimétrica. De manera ideal, el Gobierno tributaría la productividad de cada agente porque esto no generaría efectos sustitución y, por tanto, tampoco exceso de gravamen. Sin embargo, el Gobierno no puedo conocer la productividad de cada agente. Por el contrario, sí dispone de información que puede utilizar como \textit{proxy} para el conjunto heterogeneo de productividades individuales: la renta laboral de los individuos.\footnote{%
    ¿Por qué? Recordemos que, en el marco de la Síntesis Neoclásica, los modelos de Economía Laboral partían del supuesto de que, en condiciones ideales, el salario es el precio de equilibrio que iguala la utilidad marginal social que proporciona el consumo de bienes que hace posible los ingresos por renta y la productividad laboral de cada trabajador que remuneran las empresas.
}

El objetivo del modelo es determinar la estructura de los tipos y subsidios marginales que los agentes deben pagar o recibir en función de su renta observada, con el fin de hacer máximo el bienestar de la sociedad, cuyas preferencias muestran aversión a la desigualdad. Con este diseño impositivo. el gobierno pretende conseguir un determinado nivel de recaudación que financie una cuantía de gasto público considerado exógeno ($T = B$), y que por razones de simplicidad consideramos que no tiene efectos distributivos. La FBS representa una agregación del bienestar individual que pondera a los diferentes individuos de acuerdo con su posición social y donde todos los individuos tienen las mismas preferencias, aunque su productividad del trabajo sea distinta respecto de su esfuerzo. En su formulación básica. el modelo propuesto por Mirrlecs se corresponde con el siguiente programa de optimización:

\eqblock{
\begin{aligned}
    &\max_T W=\int_{w_0}^{w_m}G\left(V(w,T(w\cdot L^*))\right)f(w)\mathrm dw\\[1em]
    &\text{s.a.}\quad 
    \left\{\begin{aligned}
        &\int_{w_0}^{w_m}T(w\cdot L^*)f(w)\mathrm w\geq B\\[0.5em]
        &V(w,T(w\cdot L^*))=U(C^*,L^*)
    \end{aligned}\right.\\[2em]
    &\text{tal que},\quad (C^*,L^*)=\arg\max[U(C,L)\,,\,C=w\cdot L-T(w\cdot L)],\quad L\geq 0
\end{aligned}
}{Planificador social: Problema de optimización. Imposición directa (esquema lineal): Modelo MIRRLESS (1971)}  

La función U(C, L). creciente y cuasicóncava. recoge las preferencias de los individuos entre todas las combinaciones posibles de consumo (C) y de trabajo-ocio (L). de las que (C', C) es la preferida bajo la restricción presupuestaria que contempla el impuesto/subsidio T(-) que grava la renta bruta total observada (w • L). donde w es el salario unitario por unidad de L y V (·) es la función de utilidad indirecta del agente.

En cuanto a la función objetivo, el gobierno aspira a maximizar el valor social de las utilidades individuales mediante la FBS, C (v(w, T(w • C))), que incorpora el instrumento rcdistributivo T(w • L). La cuasi-concavidad de esta FBS refleja unas preferencias sociales con aversión a la desigualdad y. por tanto, favorables a la redistribución. El grado de concavidad de G determinará la progresividad del impuesto, con un límite inferior en la denominada FBS utilitarista. de carácter lineal y. consecuentemente, sin aversión a la desigualdad, y un límite superior en la FBS Rawlsiana. de máxima concavidad y aversión a la desigualdad. Precisamente, en esta cuasiconcavidad de la FBS subyace el lrade-o.ff entre eficiencia y equidad: a mayor aversión a la desigualdad, las preferencias sociales son favorables a una mayor redistribución a través de un diseño impositivo con mayor progresividad; sin embargo, cuanto mayor sea esta. los desincentivos al trabajo también lo son y, consecuentemente, los costes de eficiencia.\footnote{%
    En la formulación original del modelo de Mirrlees (1971 ), la relación entre esfuerzo y producto presenta rendimientos constantes a escala, lo que supone que la elasticidad de sustitución entre trabajadores es infinita y, consecuentemente, que los salarios no dependen elimina cualquier efecto redistributivo del comportamiento.
}

La primera restricción del programa sirve para fijar el nivel recaudatorio del impuesto que, como se ha apuntado, sirve para financiar el gasto público elegido por la sociedad. La segunda restricción es de incentivos. y es consecuencia del entorno de información asimétrica que plantea MÍITlees, en el que el gobierno observa y, por tanto, puede gravar las rentas percibidas por los individuos (w • L), pero no su productividad. En consecuencia, a medida que el diseüo del impuesto eleve su progresividad. los individuos se verán incentivados a comportarse como si füesen menos productivos, a trabajar menos horas. Conio se ha dicho, si la productividad fuese observable. se evitaría cualquier comportamiento estratégico por parte de los individuos para evitar su gravamen, lo que permitiría obtener la recaudación sin efecto sustitución y. por tanto. respetando el principio de neutralidad, como si se aplicase un impuesto de suma fija.

Una resolución bastante intuitiva del programa (39) puede obtenerse para unas preferencias cuasi! ineales en el consumo y una oferta de trabajo isoelástica. tal que:

$$U(C,L)=C-\frac{L^{1+\alpha^{-1}}}{1+\alpha^{-1}}$$

donde constante $\alpha$ es positiva, y la elasticidad-renta de oferta de trabajo nula. Con un impuesto sobre la renta salarial T; = T(w;; t, a), donde los parámetros de disefio / y a representan respectivamente los tipos impositivos y un subsidio (instrumentado. por ejemplo, mediante un crédito fiscal reembolsable si a > r • w¡ • L; ), la oferta de trabajo de un individuo i con una productividad reflejada en su salario unitario w; es la solución del problema primal que maximiza:

$$L^*=w_t^\alpha(1-t'(w_i\cdot L_i^*))^{\alpha}$$

Donde $t'(w_i\cdot L_i^*)$ es la derivada de la función impositiva $T(\cdot)$, es decir, el tipo marginal aplicable. A partir de esta ecuación, la resolución del problema de maximización del bienestar social anterior permite determinar el tipo marginal efectivo óptimo del impuesto ($t'^{*}$), con el que se alcanzaría el nivel de redistribución óptima:\footnote{%
    La obtención de esta solución puede verse en Atkinson y Stiglitz ( 1 980) o Diamond ( 1 998).
}

$$\frac{t'^*(w_i\cdot L_i^*)}{1-t'^*(w_i\cdot L_i^*)}=(1+\alpha^{-1})\cdot \frac{1-\int_{w_0}^{w_m}f(w)\mathrm dw}{w_i\cdot f(w)}\cdot \frac{1-s(w)}{s(w_0)} (42)$$

Donde $s(w)$ es valor marginal social promedio de la utilidad correspondiente a los individuos con una productividad superior a $w$ y $S(w_0)$ el promedio para toda la población. tal que:

$$s(w)=\frac{\int_{w_0}^{w_m}G(V(w,\,T(w\cdot L^*)))f(w)\mathrm dw}{1-\int_{w_0}^{w_m}f(w)\mathrm dw}$$

Para un impuesto sobre la renta de tipo lineal, con un único tipo marginal $t$ y una deducción reintegrable a $\alpha$ tal que $T_i=t\cdot w_i\cdot L_i^*-\alpha$, la siguiente forma simplificada permite incorporar a la expresión (42) la elasticidad-salario de la oferta de tTabajo compensada $(\varepsilon_{L,w_i}^C)$, donde $\delta$ es un término positivo con valor creciente con $\bar T$:

$$\frac{t^*}{1-t^*}=\delta \frac{\text{Cov}(S(w_0),w\cdot L^*)}{\int_{w_0}^{w_m}\varepsilon_{L,w}^{C}\cdot w\cdot L^*\cdot f(w)\mathrm dw}$$

En la expresión anterior podemos identificar, con claridad, los elementos determinantes del trade-off entre eficiencia y equidad presentes en el problema de la imposición óptima sobre la renta. El numerador del segundo término recoge el elemento distributivo del problema, y el denominador el elemento de eficiencia. Así. en el numerador. una sociedad aversa a la desigualdad dará una ponderación mayor a los individuos con salarios más bajos, de forma que la covarianza entre el valor social marginal de las utilidades y los salarios será negativa y más grande. en valor absoluto. a medida que la aversión a la desigualdad sea mayor. Por su parte. el denominador recoge una media ponderada de las elasticidades de oforta compensada de trabajo, determinante de la cuantía de los costes de eficiencia introducidos por este impuesto distorsionante, de manera que con mayores valores de estas elasticidades mayor será el exceso de gravamen del impuesto. 

De la condición expresada en (44) se obtiene que la progresividad óptima de un impuesto lineal sobre la renta con una estructura T¡ = r • w¡ • Lí - a exige fij ar un tipo marginal óptimo (r' ) mayor: a) cuanto mayor sea la aversión a la desigualdad (-G"(-)/G'(-) ): b) mayor sea la recaudación que se desea obtener (f ); y c) menor sea la elasticidad-salario promedio de la oferta de trabajo compensada (tfw ) ·

Para establecer la progresividad óptima de estas estructuras de gravamen con varios tipos marginales, definimos un vector de tipos impositivos crecientes con la renta de acuerdo con las condiciones de progresividad local propuestas en Musgrave y Thin (1948), bajo el supuesto simplificador de que cada individuo perceptor de rentas salariales será gravado con un tipo distinto $t_i$. 

Si reformulamos el programa de optimización (38). incluyendo el vector de tipos 4', tal que,
sujeto a: T(L, w, 4') = f (45)

de su resolución se obtiene la siguiente regla de imposición óptima: 


Ante la posibilidad de gravar la renta personal mediante un número finito de tipos impositivos, la estructura de gravamen $T^*$ será óptima cuando su aplicación consigue igualar las pérdidas de utilidad social marginal del último euro adicional de recaudación para todos los individuos. aw(-)/aT; = A ar(L,w,'11)/ih; ' 'vi = {l,Z , ···' N} (46)

siendo A el multiplicador langragiano de la restricción recaudatoria, el cual puede ser interpretado como el coste social asumido cuando se aumenta la recaudación en un euro adicional. 

Desarrollando ( 46). se llega a la siguiente condición. equivalente a la ecuación ( 44) para el impuesto lineal. que permite determinar el tipo marginal óptimo que debe gravar la renta del individuo i cuando se utiliza un impuesto con más de un tipo marginal:

(47)

De nuevo, obtenemos que. con este impuesto con varios tipos marginales. el tipo marginal óptimo (rt) que debe soportar la renta del individuo i ser,1 mayor: a) cuanto mayor sea la aversión a la desigualdad de la sociedad (-C"(-)/G'(-) ); b) mayor sea la recaudación que el gobierno desea obtener ( T); y c) menor sea la elasticidad-salario de la oferta de trabajo compensada del individuo i, (Ef w .) .




\paragraph{Esquema de tributación lineal: STERN (1976)}
STERN (1987) estudió un modelo similar al de EDGEWORTH, excepto en que los individuos eligen entre renta y ocio. Para simplificar el análisis, STERN supuso que los ingresos fiscales recaudados de una persona vienen dados por

Ingresos fiscales = −α + t × Renta
(16.11)

donde α y t son números positivos. Por ejemplo, supongamos que α = 3.000 euroes y t = 0,25. Entonces una persona con una renta de 20.000 euroes tendría una obligación tributaria de 2.000 euroes (= −3.000 euroes + 0,25 × 20.000 euroes). Una persona con una renta de 6.000 euroes tendría una obligación tributaria de menos 1.500 euroes (= −3.000 euroes + 0,25 × 6.000 euroes). Tal persona recibiría una transferencia de 1.500 euroes del gobierno.

En la Figura 16.4, representamos gráficamente la Ecuación (16.11) en un diagrama con la renta medida en el eje horizontal y los ingresos fiscales en el vertical. Cuando la renta es cero, la carga tributaria es negativa: el individuo recibe una transferencia del gobierno de α euroes. Después, por cada euro de renta, el individuo debe pagar t euroes al gobierno. Así, t es el tipo impositivo marginal, la proporción de un euro adicional que debe pagarse en impuestos. Debido a que la interpretación geométrica de (16.11) es una línea recta, se denomina esquema lineal del impuesto sobre la renta. En las discusiones populares, un esquema lineal del impuesto sobre la renta suele denominarse impuesto sobre la renta de tipo único. Obsérvese que, aunque el tipo impositivo marginal de un esquema tributario lineal es constante, el esquema es progresivo en el sentido de que cuanto mayor es la renta de un individuo, mayor es la proporción de la renta pagada en impuestos. (Véase el Capítulo 14). El grado de progresividad depende de los valores precisos de α y t. Valores mayores de t están asociados con sistemas tributarios más progresivos. Sin embargo, al mismo tiempo que los valores altos de t conducen a una mayor progresividad, crean mayores excesos de gravamen. El problema de la imposición óptima sobre la renta consiste en encontrar la «mejor» combinación de α y t: los valores que maximizan el bienestar social [Ecuación (16.10)] sujetos a la restricción de que se recaude una determinada cantidad de ingresos (por encima de las transferencias requeridas).

Stern [1987] encuentra que, permitiendo una cantidad modesta de sustitución entre ocio y renta, y con unos ingresos públicos requeridos equivalentes a aproximadamente el 20 por ciento de la renta, un valor de t de aproximadamente el 19 por ciento maximiza el bienestar social. Esto es considerablemente menor que el valor del 100 por ciento implícito en el análisis de Edgeworth. Incluso efectos de incentivo bastante modestos parecen tener implicaciones importantes para los tipos impositivos marginales óptimos. Por cierto, el tipo calculado por Stern también es mucho menor que los tipos impositivos marginales reales encontrados en muchos países occidentales. Por ejemplo, bajo el impuesto federal sobre la renta personal de Estados Unidos, el tipo impositivo marginal legal más alto sobre la renta en 2008 era del 35 por ciento; en algunos momentos ha sido del 90 por ciento.

De forma más general, Stern mostró que cuanto más elástica sea la oferta de trabajo, menor será el valor óptimo de t, manteniéndose todo lo demás constante. Intuitivamente, el coste de la redistribución es el exceso de gravamen que genera. Cuanto más elástica sea la oferta de trabajo, mayor será el exceso de gravamen derivado de gravarla. [Véase la Ecuación (15.4)]. Por lo tanto, una oferta de trabajo más elástica significa un coste mayor de la redistribución, de modo que debería llevarse a cabo en menor medida.

Stern también investigó cómo afectan a los resultados las funciones alternativas de bienestar social, centrándose en el impacto de asignar distintos pesos sociales a las utilidades de los ricos y de los pobres. En la Ecuación (16.10), las preferencias más igualitarias se representan asignando a las utilidades de las personas pobres pesos mayores que a las utilidades de los ricos. Un caso extremo interesante es el criterio maximin, según el cual el único individuo que recibe algún peso en la función de bienestar social es la persona con la utilidad mínima (véase el Capítulo 12). Stern encontró que el criterio maximin exige un tipo impositivo marginal de aproximadamente el 80 por ciento. Como era de esperar, si la sociedad tiene objetivos extremadamente igualitarios, se requieren tipos impositivos elevados. Sin embargo, incluso aquí, los tipos quedan por debajo del 100 por ciento.

\paragraph{Esquema de tributación progresivo: GRUBER y SAEZ (2002)}
Una limitación del análisis de Stern es que restringe el sistema del impuesto sobre la renta a tener un único tipo impositivo marginal. Gruber y Saez (2002) investigaron un modelo más general que permitía cuatro tipos impositivos marginales. Su hallazgo más interesante es que las personas en los tramos de renta más altos deberían enfrentarse a un tipo impositivo marginal más bajo que las personas en los tramos más bajos. La intuición detrás del resultado es que, al reducir el tipo impositivo marginal para las personas con ingresos altos, se las induce a ofrecer más trabajo, y el aumento de los ingresos fiscales puede utilizarse para reducir las cargas tributarias de los individuos de bajos ingresos. Es importante señalar que, aunque los tipos impositivos marginales disminuyen con la renta, los tipos impositivos medios aumentan con la renta, de modo que el sistema tributario óptimo sigue siendo progresivo. Recientemente, un cantón (estado) de Suiza implementó realmente un sistema tributario que impone tipos impositivos marginales más bajos a quienes obtienen mayores ingresos (RABUSHKA, 2003).

Esta enumeración de resultados puede transmitir una sensación algo falsa de precisión acerca de lo que los economistas realmente saben sobre el sistema tributario óptimo. Después de todo, existen muchos juicios de valor controvertidos detrás del bienestar social aditivo que el sistema tributario óptimo busca maximizar. Además, como se explica en el Capítulo 18, existe una incertidumbre considerable acerca de las elasticidades de comportamiento que son cruciales para analizar el equilibrio entre eficiencia y equidad. No obstante, calcular los tipos impositivos óptimos bajo conjuntos alternativos de supuestos es extremadamente informativo. La literatura sobre imposición óptima revela las implicaciones de supuestos éticos y conductuales alternativos y, por lo tanto, fomenta discusiones coherentes sobre política tributaria.

La utilidad marginal individual es una función decreciente del consumo individual, debido al principio de utilidad marginal decreciente. Al reducir la renta después de impuestos, unos impuestos más altos conducen a un menor consumo individual y a una mayor utilidad marginal del consumo para cualquier individuo. Si el sistema de imposición sobre la renta es tal que deja al individuo $i$ con una utilidad marginal por euro de ingresos mayor que la del individuo $j$, entonces esta formulación sugiere que los impuestos deberían reducirse para el individuo $i$ y aumentarse para el individuo $j$. Tal cambio aumentará la renta después de impuestos del individuo $i$, aumentando su consumo y reduciendo su utilidad marginal, y reducirá la renta después de impuestos del individuo j, disminuyendo su consumo y aumentando su utilidad marginal. Estos ajustes deberían continuar hasta que las ratios $MU/MR$ de todos los individuos sean iguales a $\lambda$.

El apéndice de este capítulo muestra la formulación matemática de este análisis y concluye que el sistema óptimo de imposición sobre la renta cumple la siguiente condición:

establecer las tasas del impuesto sobre la renta entre los grupos de tal manera que MUᵢ/MRᵢ = λ

donde MU es la utilidad marginal del individuo i, MR es el ingreso marginal recaudado al gravar a ese individuo y λ es el valor de los ingresos gubernamentales adicionales (como se analizó en el contexto de la imposición óptima sobre bienes). El sistema óptimo de imposición sobre la renta es aquel en el que la utilidad marginal por dólar de ingresos recaudados se iguala entre los individuos.


A pesar de su simplicidad tanto para los contribuyentes como para la administración tributaria que tiene encargada de su aplicación. el impuesto lineal sobre la renta. con un único tipo marginal y un mínimo exento, no ha sido un modelo aplicado con generalidad en prácticamente ningún país.\footnote{%
    Baldini y Rizzo (202 1) revisan las experiencias comparadas de aplicación de modelos jlat tax en países
europeos.
} Por el contrario. desde su implantación a lo largo del siglo XX. los diseños de los impuestos sobre la renta personal han contado con estructuras de gravamen que incorporan varios tipos marginales, crecientes con la renta gravada (base liquidable). Desde la última década del siglo XX, se ha consolidado una tendencia a reducir significativamente el número de tipos marginales con los que cuentan las tarifas del impuesto. siendo actualmente bastante común la adopción de modelos de tipo dual (o semi dual). donde las rentas del trabaj o son gravadas de fonna progresiva con una tarifa de varios tipos marginales y las rentas del capital con un único tipo (o con dos o tres como máximo), significativamente más reducido que los aplicados a las rentas del trabajo.



Al considerar si pasar a un sistema fiscal más progresivo, por tanto, el gobierno necesita equilibrar el hecho de que esto llevará las utilidades marginales de los ricos y los pobres hacia la igualdad (al reducir el consumo de los ricos) con el hecho de que, al gravar más fuertemente a los ricos, estos trabajarán menos, reduciendo así el ingreso marginal recaudado mediante la tributación. Este último punto se ve reforzado por lo señalado anteriormente: la ineficiencia marginal de un impuesto aumenta con la tasa impositiva, y los ricos son los elementos más productivos de la sociedad, por lo que la reducción de la producción será la mayor. Así, aumentar los impuestos a los ricos mientras se reducen para los pobres conduce a una reducción global de la eficiencia.

Como hemos visto, las ideas sobre el exceso de gravamen se aplican no solo a la tributación de bienes, sino también a la tributación de la renta. 

Otra implicación de estas reglas sobre el exceso de gravamen es que puede haber grandes costes de eficiencia al pasar de un sistema fiscal proporcional (tasas impositivas medias iguales para todos) a uno progresivo (tasas impositivas medias más altas para los ricos). Pasar a un sistema progresivo significa estrechar la base de tributación: aquella parte del impuesto que se aplica a los ricos se recauda únicamente sobre la estrecha base de renta de los ricos. Supongamos por ahora que los pobres y los ricos tienen la misma elasticidad de generación de renta con respecto a la tasa impositiva, de modo que un aumento porcentual de la tasa impositiva reduce la renta imponible en la misma cantidad porcentual para ambos grupos. Bajo este supuesto, es más eficiente gravar a todos los individuos a una tasa igual que excluir a algunos individuos de la tributación y gravar a otros individuos a una tasa impositiva más alta para compensar los ingresos perdidos. La eficiencia perdida al gravar más fuertemente a un subconjunto de individuos es mayor que la eficiencia ganada al excluir a algunos individuos de la tributación.

Este punto se ilustra mejor mediante un ejemplo. Consideremos una sociedad de dos individuos, uno de los cuales tiene un salario por hora de 10 dólares y uno de los cuales tiene un salario por hora de 20 dólares. Para ambos individuos, un aumento del $10\%$ en los salarios los lleva a ofrecer un $10\%$ más de oferta de trabajo: su elasticidad de la oferta de trabajo con respecto a los salarios después de impuestos es $1$. La elasticidad de la demanda de trabajo con respecto a los salarios es $–1$: por cada aumento del $10\%$ en los salarios, las empresas demandan un $10\%$ menos de horas de trabajo. Gráficamente:

\imagenfit{Fig2.png}{Esquema de tributación progresiva: ¿mayores costes de eficiencia?}{}

Sin ningún impuesto, el trabajador de salario bajo trabaja $1.000$ horas ($H_1$) y gana $10€$ por hora ($W_1$), como se muestra en el panel (a), mientras que el trabajador de salario alto también trabaja $1.000$ horas ($H_1$), pero gana $20€$ por hora ($W_1$), como se muestra en el panel (b).\footnote{
Los resultados gráficos se resumen en la siguiente tabla:

\begin{center}
\begin{tabular}{|p{0.1\linewidth}p{0.09\linewidth}p{0.09\linewidth}|p{0.09\linewidth}p{0.09\linewidth}|p{0.09\linewidth}p{0.09\linewidth}|p{0.09\linewidth}|}
\hline
 &  & & \multicolumn{2}{c}{\textbf{Panel (a): Ingresos bajos}} & \multicolumn{2}{c}{\textbf{Panel (b): Ingresos altos}} &  \\
    \cline{4-7}
 & {\scriptsize\textbf{Impuesto: }$\mathbf{W<10.000}$} &{\scriptsize\textbf{Impuesto}: $\mathbf{W>10.000}$} & {\scriptsize\textbf{Horas ofertadas}} & {\scriptsize\textbf{Exceso de gravamen}} & {\scriptsize\textbf{Horas ofertadas}} & {\scriptsize\textbf{Exceso de gravamen}} & {\scriptsize\textbf{Exceso de gravamen total}} \\
 \hline
 {\scriptsize\textbf{Sin impuesto}} & $0\%$ & $0\%$ & $1.000\, (H_1)$ & $0€$ & $1.000\,(H_1)$ & $0€$ & $0€$ \\
 \hline
 {\scriptsize\textbf{Esquema lineal (proporcional)}} & $20\%$ & $20\%$ & $894\, (H_1)$ & $115,71€$ & $894\,(H_1)$ & $321,42€$ & $347,13€$ \\
 \hline
 {\scriptsize\textbf{Esquema progresivo}} & $0\%$ & $60\%$ & $1.000\, (H_1)$ & $0€$ & $837\,(H_1)$ & $566,75€$ & $566,75€$ \\
 \hline
\end{tabular}
\end{center}
}

Supongamos que el gobierno impone un impuesto sobre las nóminas del $20\%$ sobre todas las ganancias de un trabajador (un impuesto proporcional). Este impuesto reduce la renta que ambos trabajadores reciben del trabajo, haciéndolos menos dispuestos a ofrecer trabajo, como se refleja en el desplazamiento de la curva de oferta de trabajo de S1 a S2 en ambos paneles. Para el trabajador de salario bajo, la oferta de trabajo cae de 1.000 horas (H1) a 894 horas (H2), a un salario antes de impuestos más alto de 11,18 dólares. Como resultado, el impuesto ha creado una pérdida de peso muerto del área BAC, que tiene un área de 115,71 dólares. Para el trabajador de salario alto, la cantidad de trabajo ofrecida también cae de 1.000 horas a 894 horas a un salario antes de impuestos más alto de 22,36 dólares, creando una pérdida de peso muerto del área EDF, que tiene un área de 231,42 dólares.

Pista rápida ¿Por qué es mayor la pérdida de peso muerto para el trabajador de salario más alto a pesar de la misma reducción en las horas trabajadas? En un mercado de trabajo competitivo, la tasa salarial es igual al producto marginal del trabajo, por lo que el trabajador de salario alto tiene un producto marginal del trabajo más alto. Como resultado, la sociedad pierde más eficiencia cuando la trabajadora de salario alto reduce sus horas (a un producto marginal de 20 dólares por hora) que cuando la trabajadora de salario bajo reduce sus horas (a un producto marginal de 10 dólares por hora).

Ahora supongamos que la sociedad decide cambiar a un esquema fiscal progresivo: ningún impuesto sobre los primeros 10 dólares de ganancias por hora, pero una tasa impositiva del 60% sobre los siguientes 10 dólares de ganancias por hora. Este sistema fiscal progresivo recaudará la misma cantidad de ingresos que el impuesto proporcional del 20% sobre todas las ganancias. Sin embargo, estos dos sistemas fiscales diferentes tienen consecuencias de eficiencia muy diferentes.

La Figura 20-7 ilustra este punto. En el panel (a), el cambio al esquema fiscal progresivo aumenta los beneficios de trabajar para la trabajadora de salario bajo (porque ya no está gravada) y hace que su curva de oferta de trabajo se desplace de nuevo a S1. Ahora ha vuelto a su óptimo original en el punto A: 1.000 horas de trabajo a 10 dólares por hora. La pérdida de peso muerto asociada con la perceptora de salario bajo cae de 115,71 dólares a cero. En el panel (b), el cambio en el esquema fiscal reduce aún más los beneficios de trabajar para la trabajadora de salario alto y desplaza su curva de oferta hacia dentro hasta S3, de modo que sus horas caen de 894 a 837 (H3). Este cambio aumenta el tamaño del triángulo de pérdida de peso muerto a GDI, añadiendo el área del trapecio GEFI. La DWL en el panel (b) tiene ahora un área total de 566,75 dólares.

Este cambio en el esquema fiscal ha reducido la DWL en 115,71 dólares para la trabajadora de salario bajo y la ha aumentado en 335,33 dólares para la trabajadora de salario alto. En términos netos, la DWL ha aumentado de 347,13 dólares a 566,75 dólares, un aumento del 63% (219,62 dólares).

El gran aumento de la DWL surge porque este impuesto más progresivo se recauda sobre una base imponible más pequeña. El impuesto proporcional se recauda sobre un total de 30.000 dólares de ganancias (10.000 dólares de ganancias para el trabajador de salario bajo y 20.000 dólares de ganancias para el trabajador de salario alto). El impuesto progresivo se recauda sobre un total de solo 10.000 dólares de ganancias (los segundos 10.000 dólares de ganancias del trabajador de salario alto). Para recaudar la misma cantidad de ingresos sobre esta base imponible más pequeña, el impuesto progresivo debe imponer una tasa impositiva más alta, y una tasa impositiva más alta significa una pérdida de peso muerto marginal más alta. La contribuyente de salario bajo ve una reducción en su tasa impositiva marginal del 20% a cero, mientras que la contribuyente de salario alto ve un aumento en su tasa impositiva marginal del 20% al 60%. Debido a que la pérdida de peso muerto aumenta con el cuadrado de la tasa impositiva, las ganancias de eficiencia derivadas de reducir la tasa impositiva del salario bajo en un 20% son menores que los costes de eficiencia de aumentar la tasa impositiva del salario alto en un 40%.

Este ejemplo ilustra el punto más amplio de que cuanto más cargas los impuestos sobre una fuente, más rápidamente aumenta la pérdida de peso muerto. Según esa lógica, los sistemas fiscales más eficientes son aquellos que distribuyen la carga de la tributación de la manera más amplia posible, de modo que la tasa impositiva, el impulsor de la pérdida de peso muerto, pueda minimizarse. El principio rector para una tributación eficiente es crear un campo de juego amplio y nivelado en lugar de gravar a algunos grupos o bienes de manera particularmente elevada y a otros no gravarlos en absoluto.


\textcolor{red}{\textbf{OJO-> }Ejemplo más sencillo (GUBER)}
Imagina un mundo sin tributación, en el que el gobierno quiere introducir un pequeño impuesto sobre la renta del 1%. En un mundo así, MU/MR es mucho menor para una persona rica que para una pobre. MU es mucho menor para la persona rica porque él o ella ya tiene un nivel de consumo muy alto, y MR también es mayor para esa persona rica porque un impuesto del 1% recauda más dinero sobre una base de renta más alta. Por tanto, la ratio MU/MR es mucho menor para los ricos que para los pobres.

Sin embargo, a medida que gravamos más a los ricos, su ratio de MU/MR aumenta. El numerador aumenta porque su consumo está cayendo, por lo que MU es mayor. El denominador cae a través de las respuestas conductuales de la oferta de trabajo: el impuesto más alto hace que trabajen menos intensamente, la base imponible se hace más pequeña y, por tanto, el impuesto recauda menos ingresos. En algún momento, cuando el impuesto sobre los ricos sea suficientemente alto, su MU/MR caerá realmente por debajo del de los pobres. Las tasas impositivas deberían ser más altas para los ricos que para los pobres, pero el gobierno no debería «exprimir» totalmente a los ricos; debería elevar su MU, pero no tanto como para que su MR procedente del impuesto se vuelva muy pequeño (o incluso negativo, como en el lado equivocado de la curva de Laffer).

\imagenfit{Fig3.png}{Imposición óptima: Esquema progresivo}{GRUBER}

La Figura 20-10 ilustra este punto. En el eje x de este diagrama está la tasa impositiva, y en el eje y está la ratio de MU/MR. Las curvas del gráfico muestran cómo cambia MU/MR a medida que aumentan las tasas impositivas. Estas curvas tienen pendiente ascendente. A medida que aumentan las tasas impositivas, MU aumenta (como resultado de la utilidad marginal decreciente), y MR cae (como resultado de las reducciones de la oferta de trabajo cuando aumentan las tasas impositivas). De hecho, el aumento simultáneo del numerador y la reducción del denominador conducen a curvas que no solo tienen pendiente ascendente, sino que tienen pendiente ascendente a una tasa creciente (las curvas se vuelven más pronunciadas a medida que aumentan las tasas impositivas). En el límite, con una tasa impositiva del 100%, MU/MR es infinito para todos (la curva es perfectamente vertical) porque el impuesto no recauda dinero.

Consideremos dos individuos, el Sr. Rico, que tiene una renta alta, y la Sra. Pobre, que tiene una renta baja. La curva del Sr. Rico comienza por debajo de la de la Sra. Pobre, porque el Sr. Rico tiene una MU menor por cada dólar adicional de renta. Si gravamos a ambos individuos al 10%, entonces MU/MR para el Sr. Rico está muy por debajo de MU/MR para la Sra. Pobre.

A medida que aumentamos el impuesto sobre el Sr. Rico, su MU aumenta y su MR cae. Para cuando la tasa impositiva sobre el Sr. Rico es del 20%, su MU/MR es igual al de la Sra. Pobre a la tasa del 10%. Este es el par óptimo de tasas del impuesto sobre la renta porque MU/MR = λ para ambos contribuyentes.

\subsubsection{Impuesto sobre el gasto personal: Desarrollo analítico (CHAMLEY, 1986; JUDD (1985))}
Aunque el impuesto sobre la renta personal ha sido una herramienta tradicional para recaudar ingresos y redistribuir la riqueza, se han analizado también otras alternativas, como las reglas de imposición óptima sobre el gasto personal. Este tipo de tributación se sustenta sobre la base de los esquemas de imposición progresiva de la renta y ha sido valorado en diversos informes de Política Fiscal muy destacados, como los Informes de las comisiones Meade (1978) para el Reino Unido, Lodin (1978) para Suecia o el Informe Bradford ( 1986) para Estados Unidos. 

Aunque no se aplica en ningún país, la utilización del gasto como alternativa a la renta personal en el diseño del impuesto supone valorar las implicaciones sobre el bienestar social que tendría la utilización del consumo presente como magnitud para evaluar la capacidad de pago de los individuos.  La defensa teórica de la imposición óptima sobre el gasto/consumo frente a la imposición sobre la renta parte de una idea bastante distinta (aunque conectada) con MIRRLEES: \textbf{gravar la renta no solo grava el trabajo, sino también el rendimiento de ahorrar, y por tanto puede distorsionar innecesariamente la decisión intertemporal entre consumir hoy y consumir mañana}.

La intuición central puede verse así. Un individuo obtiene renta laboral wL y decide cuánto consumir hoy y cuánto ahorrar:

wL = C_1 + S

y el ahorro permite consumir mañana:

C_2=(1+r)S.

Si introduces un impuesto sobre la renta que grava también los rendimientos del capital a un tipo \tau:

C_2=[1+r(1-\tau)]S.

Por tanto, el impuesto sobre la renta reduce el precio relativo del consumo presente respecto al consumo futuro. No solo distorsiona la elección trabajo–ocio —que es la preocupación de Mirrlees—, sino también la elección consumo presente–consumo futuro.

En cambio, con un impuesto uniforme sobre el consumo, simplificando,

(1+t)C_1 + \frac{(1+t)C_2}{1+r}=wL,

el rendimiento intertemporal sigue siendo 1+r. Consumir mañana no queda penalizado relativamente respecto a consumir hoy. El impuesto reduce el poder adquisitivo, pero no introduce una cuña adicional entre consumo presente y futuro.

Ahí está el primer gran argumento.

1. El famoso resultado de Chamley–Judd

La versión más potente de esta intuición aparece en los modelos de imposición dinámica óptima de Chamley (1986) y Judd (1985).

En determinados modelos, el resultado es:

\boxed{\tau_K^* \rightarrow 0}

en el largo plazo.

¿Por qué? Porque un impuesto permanente sobre el rendimiento del capital genera una distorsión que se acumula con el horizonte temporal.

Imagina que tienes €1 y puedes consumirlo hoy o invertirlo durante n períodos. Sin impuesto:

C_n=(1+r)^n.

Con impuesto sobre el rendimiento:

C_n=[1+r(1-\tau_K)]^n.

La cuña entre ambos crece con n. Cuanto más lejano sea el consumo, mayor es la penalización fiscal implícita.

Por eso, desde esta perspectiva, gravar permanentemente el rendimiento del ahorro es una forma especialmente costosa de recaudar.

2. Atkinson–Stiglitz: otra vía para llegar a una conclusión parecida

Hay otro resultado fundamental, el teorema de Atkinson–Stiglitz (1976).

Supón que tienes individuos con diferentes productividades w, que el gobierno no observa. Eso nos devuelve exactamente al mundo de Mirrlees.

El gobierno puede utilizar:

* un impuesto no lineal sobre la renta laboral T(wL);
* impuestos sobre distintos bienes de consumo.

Atkinson y Stiglitz muestran que, bajo ciertas condiciones —especialmente separabilidad entre ocio y consumo y preferencias suficientemente homogéneas—, si el gobierno ya dispone de un impuesto óptimo no lineal sobre la renta laboral, no obtiene nada gravando diferencialmente los distintos bienes de consumo.

Esquemáticamente, si:

U=u(C_1,C_2)-v(L),

entonces, bajo sus supuestos,

\boxed{\text{imposición diferencial óptima sobre bienes}=0}.

Y si interpretas C_1 y C_2 como consumo en distintos momentos, gravar el ahorro equivale a gravar más el consumo futuro que el presente.

De ahí sale una conclusión muy importante:

\boxed{\tau_K=0}

bajo esas condiciones.

Fíjate en que Chamley–Judd y Atkinson–Stiglitz llegan a resultados relacionados por mecanismos teóricos diferentes.

3. Pero entonces, ¿por qué no eliminar simplemente todos los impuestos sobre la renta y poner un IVA enorme?

Aquí está el matiz fundamental. Estos resultados no implican necesariamente que no debamos gravar la renta laboral.

Lo que atacan especialmente es la imposición sobre la renta del capital.

Puedes pensar:

Y = wL + rK.

Un impuesto sobre la renta comprensiva grava:

T=T(wL+rK).

Mientras que buena parte de la literatura de imposición óptima llega, bajo determinados supuestos, a algo conceptualmente más cercano a:

\boxed{\text{gravar } wL,\qquad \text{no gravar } rK.}

Y un impuesto sobre el gasto puede conseguir algo parecido.

Si defines:

\text{Gasto}= \text{Renta}-\text{Ahorro},

un impuesto personal sobre el gasto permite deducir el ahorro cuando se realiza y grava los recursos cuando finalmente se destinan al consumo.

Por eso el debate teórico relevante no es simplemente:

\text{Income tax vs. VAT}

sino algo más profundo:

\boxed{\text{¿Debemos gravar la renta cuando se genera o cuando se consume?}}

4. Hay además un argumento de neutralidad intertemporal

Supón dos personas idénticas que ganan €100.000.

Ana consume inmediatamente sus €100.000. Carlos consume €50.000 y ahorra €50.000 para consumirlos dentro de 20 años.

Con un impuesto sobre la renta del capital, Carlos termina pagando más impuestos simplemente porque pospuso su consumo.

Desde la perspectiva de ciertos modelos de imposición óptima, esto es difícil de justificar: ambos obtuvieron inicialmente los mismos recursos, pero el sistema penaliza al que decide trasladar consumo hacia el futuro.

Un impuesto sobre el gasto intenta ser neutral respecto a esa decisión:

C_{\text{hoy}}\quad\leftrightarrow\quad C_{\text{mañana}}.

La conexión con lo que estábamos viendo de Mirrlees/Stern

Aquí creo que está la clave para encajar toda esta literatura.

Mirrlees/Stern parten de:

\text{redistribución}
\quad\leftrightarrow\quad
\text{distorsión de la oferta laboral}.

La literatura posterior pregunta: si necesitamos redistribuir, ¿qué bases imponibles debemos utilizar para hacerlo con el menor coste de eficiencia posible?

Y algunos resultados clásicos responden, bajo supuestos fuertes:

\text{Grava la capacidad/renta laboral de forma óptima}

pero

\boxed{\text{no añadas una distorsión sobre la decisión ahorro-consumo si no te aporta información redistributiva adicional}.}

Ese último inciso —“si no aporta información adicional”— es crucial. Es precisamente por donde la literatura más moderna ha cuestionado los resultados de impuesto cero al capital: si el patrimonio, el ahorro o los rendimientos del capital contienen información sobre la capacidad de generar renta que el gobierno no observa, si existen herencias, rentas económicas, heterogeneidad de rendimientos, restricciones de liquidez, etc., entonces gravar el capital puede volver a ser óptimo.

Así que la defensa académica del expenditure tax no es simplemente “el consumo es una base mejor que la renta”. El argumento teórico profundo es: si puedes conseguir tus objetivos redistributivos gravando trabajo/consumo sin distorsionar además el margen intertemporal, no hay razón para introducir esa segunda distorsión mediante un impuesto sobre el rendimiento normal del ahorro.



\subsection{Teoría de la Imposición Óptima: Teorema de ATKINSON y STIGLITZ (1976)}
Antes de pasar a ver otras cuestiones de la Teoría de la Tributación Óptima, que poco tienen que ver con el establecimiento de reglas de tributación directa y/o indirecta, conviene presentar uno de los resultados más interesantes sobre estas dos grandes categorías de figuras impositivas. Se trata del Teorema de ATKINSON y STIGLITZ (1976) y constituye una aportación muy relevante para la consecución de los objetivos de eficiencia y equidad, mediante la combinación óptima de estas dos categorías. 

En una economía en la que los individuos solamente se diferencian por su capacidad de obtención de rentas, y bajo el supuesto de que la utilidad es separable en la cantidad de trabajo ofertado y en el consumo de bienes, los autores demuestran que: si las diferencias entre individuos pueden tratarse óptimamente mediante un impuesto no lineal sobre la renta laboral, la composición del consumo no aporta información adicional sobre esas diferencias (separabilidad) y, por lo tanto, no existe una justificación redistributiva para gravar diferencialmente los bienes.\footnote{%
    Fijémonos en que la separabilidad es lo que hace funcionar el resultado. Si, por ejemplo, las personas de alta capacidad tienen sistemáticamente preferencias distintas por determinados bienes, o ciertos bienes son complementarios del ocio o del trabajo, el consumo sí puede revelar información útil sobre la capacidad o ayudar a corregir la distorsión laboral. En ese caso, ATKINSON-STIGLITZ deja de aplicarse y los impuestos indirectos diferenciados pueden volver a ser óptimos.
}

¿Por qué desaparece entonces la utilidad redistributiva de los impuestos indirectos diferenciados? Imaginemos dos tipos de individuos: $w_H>w_L$. El gobierno quiere redistribuir de $H$ hacia $L$. Para ello dispone de un impuesto no lineal sobre la renta: $T(y)$. Supongamos ahora que además piensa: los ricos consumen más del bien $X$, así que voy a gravar $X$ más fuertemente para redistribuir. ATKINSON-STIGLITZ demuestran que, bajo sus supuestos, esto no mejora el resultado. ¿Por qué? Porque el gobierno ya puede realizar la redistribución deseada mediante $T(y)$ y, si introduce además $t_X>t_Y$, está distorsionando la elección de consumo sin obtener una ventaja adicional en términos de redistribución/incentivos. Por tanto, en el óptimo: $t_1=t_2=\cdots=t_n$

En definitiva, una vez que sabemos cuánto ingreso obtiene una persona y diseñamos óptimamente $T(y)$, saber si gasta sus $1.000€$ en restaurantes, libros, ropa o muebles no nos ayuda a distinguir mejor entre individuos de alta y baja capacidad. Por tanto, gravar esos bienes de manera diferente añade una distorsión sin mejorar el problema informativo que limita la redistribución.\footnote{%
    Además, si comparamos la capacidad redistributiva del impuesto sobre la renta con la del impuesto sobre el consumo, para que la de este fuese superior a la del primero, el impuesto sobre el consumo debería gravar con tipos mayores los bienes con una elasticidad-renta mayor. Sin embargo, de acuerdo con la evidencia empírica, conocemos que estos bienes suelen presentar elasticidades-precio mayores, por lo que las distorsiones que generaría su gravamen serían, comparativamente, también más elevadas. Por consiguiente, con un adecuado diseño óptimo del impuesto sobre la renta, las ganancias de bienestar social netas aportadas por la utilización de la imposición indirecta con fines redistributivos sería pequeña o posiblemente negativa.
}

No obstante, con posterioridad a este teorema, la literatura muestra que, si se introducen otros supuestos razonables en un contexto real de aplicación de los sistemas fiscales (evasión fiscal, heterogeneidad en otros factores, etc), pueden encontrarse argumentos sólidos para la coexistencia de imposición directa progresiva e imposición indirecta diferenciada.\footnote{%
    En este sentido, como apuntan VÁZQUEZ, VULOVIC y LIU (2010), debe asumirse que el diseño de sistemas fiscales óptimos requiere contar con impuestos directos e indirectos, aunque no existen reglas concluyentes que permitan establecer cuál debe ser la combinación óptima de ambas formas.
}



\clearpage
\section{Sistemas Fiscales: Otras consideraciones}
\puntosclave{
    La teoría de la evasión fiscal óptima busca explicar cómo los individuos toman decisiones sobre si evadir o no impuestos y cómo el diseño del sistema tributario y de la administración encargada de su gestión afecta a sus incentivos para hacerlo.
    
    Un contribuyente que considera la posibilidad de incumplir sus obligaciones fiscales decide qué cantidad de renta le interesa evadir maximizando la utilidad esperada de su renta disponible. de acuerdo con unas preferencias con aversión al riesgo.
    
    La resolución del problema de evasión fiscal óptima permite deducir algunas relaciones de interés para el diseño de la política de administración tributaria: a) al aumentar el porcentaje de sanción la cantidad evadida óptima disminuye; b) el aumento de la probabilidad de ser inspeccionado hace que la renta evadida óptima disminuya; c) al incrementarse la renta. la cantidad evadida óptima aumenta; d) cuando la sanción se establece sobre la cantidad de renta evadida, el aumento del tipo impositivo tienen un efecto ambiguo sobre la evasión, mientrns que si se aplica sobre la cuota no satisfecha, un incremento del tipo impositivo hace disminuir el importe óptimo de renta evadida.

    En un modelo de administración tributaria óptima que incorpora la intensidad de control al problema de elección del tipo impositivo óptimo. bajo los supuestos de que gravar las bases impositivas es una actividad costosa en presencia de comportamientos evasores y que los costes de eficiencia aumentan de fo1ma no lineal al hacerlo el tipo impositivo, se obtienen las siguientes reglas de diseño óptimo: a) para una aplicación impositiva óptima. es condición suficiente. aunque no necesaria. que la intensidad de control y el tipo impositivo estén relacionados tal que. al aumentar la primera, aumente el segundo, lo que hará disminuir la renta evadida; b) en el margen, es mejor obtener ingresos impositivos adicionales, netos de costes de administración. incrementando el esfuerzo de control por parte de la administración tributaria. que aumentando el tipo impositivo; c) el nivel óptimo de control tributario está condicionado por las posibles relaciones de complementariedad o sustituibilidad entre la evasión y la elusión que presente el diseño impositivo.
}


La imposición óptima es una teoría puramente normativa. No pretende predecir cómo son los sistemas fiscales del mundo real ni explicar cómo surgen estos sistemas fiscales. La teoría presta poca atención al entorno institucional y político en el que se elabora la política fiscal. Holcombe [2002] sostiene que, en presencia de instituciones políticas del mundo real, las recomendaciones de política basadas en la lógica de la imposición óptima pueden en realidad reducir el bienestar.

Supongamos que, en una determinada sociedad, hay tres mercancías, X, Y y ocio. La oferta de trabajo es totalmente fija y, por lo tanto, la renta es fija. Actualmente, esta sociedad aplica un impuesto sobre X, pero su constitución prohíbe gravar Y. Al observar esta situación, un estudiante de la teoría de la imposición óptima podría decir algo como: «Están aplicando un sistema fiscal ineficiente. Debido a que la oferta de trabajo es totalmente fija, podrían no tener ningún exceso de gravamen si gravaran X e Y a tipos iguales: un impuesto sobre la renta. Recomiendo que reduzcan el impuesto sobre X e impongan un impuesto al mismo tipo sobre Y. Fijen los tipos de manera que se recaude la misma cantidad de ingresos que antes».

Supongamos, sin embargo, que los ciudadanos sospechan que, si permiten la imposición sobre Y, sus políticos no reducirán el tipo impositivo sobre X. Más bien, simplemente aprovecharán la oportunidad de gravar algo nuevo para hacer que los ingresos fiscales sean lo más grandes posible. Como vimos en el Capítulo 6, ciertas teorías del sector público sugieren que quienes dirigen el gobierno pueden maximizar y maximizarán los ingresos fiscales a pesar de los deseos de la ciudadanía. Por lo tanto, al impedir constitucionalmente la imposición sobre Y, los ciudadanos pueden estar protegiéndose racionalmente contra un sector público ineficientemente grande. En otras palabras, si los ciudadanos no confían en el gobierno, lo que parece ineficiente a través de la lente de la imposición óptima sobre las mercancías puede ser eficiente en un contexto más amplio.⁸ De hecho, existe cierta evidencia de que los gobiernos con sistemas fiscales que generan grandes excesos de gravamen tienden a crecer más lentamente que los gobiernos con sistemas fiscales eficientes [Becker y Mulligan, 2003].

Las cuestiones relacionadas con estas consideraciones pueden ayudar a explicar, en parte, la controversia actual sobre el tratamiento fiscal de las compras realizadas en Internet. Los defensores de la imposición sobre Internet sostienen que un bien comprado en una tienda es esencialmente la misma mercancía que el mismo bien comprado en Internet. Gravar el primero pero no el segundo distorsiona las elecciones de los consumidores entre las dos modalidades de compra y, por lo tanto, crea un exceso de gravamen. Los opositores sostienen que gravar las ventas por Internet simplemente alimentaría aumentos en el tamaño del sector público, que ya es ineficientemente grande.

Esta discusión está relacionada con la inconsistencia temporal de la política óptima, que se produce cuando el gobierno no puede implementar una política fiscal óptima porque la política declarada es inconsistente con los incentivos del gobierno a lo largo del tiempo. Considérese una propuesta realizada por el gobierno de Colombia en 2002. Para sofocar una rebelión, se aplicaría un impuesto del 1,2 por ciento del valor de su capital a todos los individuos y empresas cuyos activos superaran el equivalente a 60.000 dólares. Es importante señalar que el impuesto debía imponerse una sola vez; no se repetiría en el futuro. Aunque presumiblemente los capitalistas no estarían satisfechos de pagar el impuesto, parecería no tener ningún impacto sobre sus incentivos actuales para ahorrar para el futuro. Un impuesto de este tipo es, en efecto, un gravamen de suma fija y, por lo tanto, plenamente eficiente.

Sin embargo, existe un problema. El gobierno colombiano tiene un incentivo para incumplir su promesa de que el impuesto solo se aplicaría una vez y hacer exactamente el mismo truco el año siguiente, recaudando todavía más ingresos sin un exceso de gravamen. Así, la política fiscal declarada es inconsistente con los incentivos del gobierno a lo largo del tiempo. Peor aún, los capitalistas se dan cuenta de que el gobierno tiene un incentivo para incumplir su promesa. Cambiarán su comportamiento de ahorro para reflejar la expectativa de que cuanto más ahorren ahora, más impuestos se les cobrarán el próximo año. Debido a que el impuesto esperado cambia el comportamiento, introduce una ineficiencia.

En resumen, a menos que el gobierno pueda prometer de manera creíble que no incumplirá su promesa, no puede llevar a cabo la política fiscal plenamente eficiente. Para evitar este problema de inconsistencia temporal, el gobierno debe ser capaz de comprometerse a comportarse de determinadas maneras en el futuro. ¿Cómo puede hacerse esto? Un posible enfoque consiste en promulgar disposiciones constitucionales que prohíban al gobierno incumplir sus promesas. Sin embargo, mientras el gobierno tenga un incentivo subyacente para incumplirlas, las sospechas permanecerán, frustrando los intentos de llevar a cabo una política eficiente. Estas consideraciones sugieren que debe tenerse en cuenta la credibilidad del sistema político antes de formular recomendaciones basadas en la teoría de la imposición óptima.




Tras la extensión de la TIO a la imposición distribución con metas distributivas a comienzos de la década de los setenta del siglo pasado, su marco de análisis ha ido ampliándose para cubrir diferentes ámbitos del diseño impositivo. Esta ampliación de sus áreas temáticas ha atendido a la incorporación del principio impositivo de sencillez de gestión, cubriendo dentro del mismo tanto la faceta relacionada con el cumplimiento de los contribuyentes, como la relacionada con la administración del sistema tributario.


\subsection{Coste Marginal de los Fondos Públicos}

Una suposición implícita en los modelos que hemos estado estudiando es que recaudar impuestos no implica costes. Esto es claramente falso. Las autoridades fiscales necesitan recursos para hacer su trabajo. Los contribuyentes también incurren en costes, incluidos los desembolsos para contables y abogados fiscalistas, así como el valor del tiempo dedicado a rellenar declaraciones de impuestos y llevar registros.

Los costes de administrar el impuesto sobre la renta en Estados Unidos son bastante bajos. Por ejemplo, el Servicio de Impuestos Internos gasta solo unos 44 centavos para recaudar cada 100 dólares en impuestos. Sin embargo, los costes de cumplimiento de la imposición sobre la renta personal son bastante sustanciales. Estos costes de cumplimiento incluyen el tiempo dedicado a la preparación de los impuestos y el coste de elementos tales como el asesoramiento profesional y los manuales de preparación. La evidencia procedente de encuestas sugiere que el coste total de cumplimiento del impuesto sobre la renta es aproximadamente el 10 por ciento de los ingresos [Kaplow, 2008a], o unos 122.000 millones de dólares en 2008.

Claramente, la elección de los sistemas de impuestos y subsidios debería tener en cuenta los costes administrativos y de cumplimiento. Incluso los sistemas que parecen equitativos y eficientes (en el sentido del exceso de gravamen) podrían ser indeseables porque son excesivamente complicados y costosos de administrar. Considérese la posibilidad de gravar la producción doméstica —limpieza del hogar, cuidado de los niños, etcétera—. Como se sugirió en el Capítulo 15, el hecho de que el trabajo de mercado esté gravado pero el trabajo doméstico no lo esté crea una distorsión considerable en la asignación del trabajo. Además, gravar de manera diferente sobre la base de la elección del lugar de trabajo viola algunas nociones de equidad horizontal. No obstante, las dificultades implicadas en valorar la producción doméstica crearían unos costes administrativos tan enormes que la idea es inviable.

Desafortunadamente, los problemas administrativos a menudo reciben una atención insuficiente. Un caso clásico fue el impuesto federal sobre bienes de lujo aplicado a las joyas nuevas, promulgado en 1990. El impuesto se aplicaba únicamente a la parte del precio que excedía los 10.000 dólares, y solo los artículos usados como adorno estaban sujetos al impuesto. Como señaló un comentarista, el impuesto era una pesadilla administrativa: «las gemas sueltas y las reparaciones no están gravadas; el valor de mercado después de una modificación importante sí lo está. Así, […] puede que se le cobre un impuesto si hace colocar en una nueva montura las gemas del broche de su abuela. Pero no se le cobrará si reemplaza un diamante de 30.000 dólares perdido de un anillo; eso es una reparación».¹⁰ ¡Los costes para el Servicio de Impuestos Internos de recaudar el impuesto sobre bienes de lujo pueden haber superado los ingresos recaudados! El impuesto fue finalmente derogado en 1993.

Obviamente, ningún sistema fiscal carece de costes de administración; el truco consiste en encontrar la mejor disyuntiva entre el exceso de gravamen y los costes administrativos. Por ejemplo, administrar un sistema de impuesto sobre las ventas en el que cada mercancía tenga su propio tipo podría ser muy engorroso, a pesar de que este es el enfoque general prescrito por la regla de Ramsey. Cualquier reducción del exceso de gravamen que surja de diferenciar los tipos impositivos debe compararse con los costes administrativos incrementales.


\textcolor{red}{\textbf{Origen:} ICEX-CECO}
El concepto de coste marginal de los fondos públicos (CMFP) es una herramienta esencial en el análisis impositivo y se ha desarrollado y refinado a lo largo del tiempo. El CMFP se utiliza para evaluar comparativamente el impacto en el bienestar generado por la recaudación de impuestos, a través de la relación entre la recaudación obtenida (7) y el exceso de gravamen (EG) originado. De forma intuitiva. el CMFP mide, en términos monetarios. el coste económico de recaudar un euro adicional usando un determinado impuesto. interpretado como el coste de opo11unidad de utilizar ese euro como recurso para financiar el gasto público en lugar de su libre disposición por los agentes privados:
CMFP =��T (17)

El concepto seminal de esta medida fue introducido por Browning ( 1976, 1987). definida en un marco de análisis de '·reforma diferencial"', en el que se busca calcular el coste de eficiencia que podría evitarse reformando el sistema fiscal. manteniendo el nivel de recaudación total y la composición del gasto público financiado con la misma. Por tanto. el punto ele partida es la restricción presupuestaria que iguala los ingresos impositivos ( 7). obtenidos gravando el consumo de bienes o la renta, y nivel de gasto público en bienes y servicios públicos a financiar (G):
( 1 8)

Sea W la función de bienestar social que representa las preferencias del conjunto de la sociedad, expresada en términos del consumo de sus n miembros.

$W=W(C_1,C_2,\dots,C_{n-1},C_n)$

el CMFP se puede calcular como el cambio en bienestar resultante de incrementar en una unidad los ingresos públicos elevando la recaudación. tal que:

CMFP = aw /ar (20)

Incorporando la restricción presupuestaria a través del multiplicador de Lagrange (il), se obtiene la función lagrangiana.
L = W + i\(T - C) ( 2 1 )

que diferenciamos respecto del consumo individual (C¡ ) y de los impuestos (T), obteniendo las siguientes condiciones de optimización:
a1, _ aw i\ ac
ac; - ac; - ac, (22)
(23)

Resolviendo para i\. se obtiene el tipo impositivo óptimo (T"') que maximiza W, cumpliendo con la restricción de financiación del gasto ( 1 8). 

(24) 

de donde se obtiene que. para el individuo i. el coste marginal de financiar el gasto público es: 

CMFP- = aw t 8 C¡ (25) 

De esta forma. la expresión (25) proporciona una medida de la pérdida de bienestar experimentada por el individuo i como consecuencia destinar un último euro adicional a la financiación del gasto público. Esta pérdida de bienestar es consecuencia del cambio en su nivel de forzado por el pago del impuesto. 

Desde su aparición. han ido surgiendo medidas alternativas del concepto de CMFP, en gran medida como consecuencia de las necesidades planteadas por el análisis empírico y las limitaciones de los métodos analíticos.\footnote{%
    En González-Páramo (2003) se ofrece una completa revisión de las alternativas disponibles en la literatura, con sus potencialidades y limitaciones. Recientemente. Bastani (2023) ha publicado una interesante guía actualizada. centrada fundamentalmente en aclarar las d iscrepancias sobre su definición y medida, tanto desde el punto de vista teórico como empírico.
}

Como señala González-Páramo (2003), a nivel agregado y para una estructura impositiva determinada, un CMFP correspondiente al sistema fiscal que fuese mayor que el beneficio marginal del gasto público estaría revelando que la dimensión del sector público en términos de gasto es excesiva. A nivel micro, fijada una opción de financiación de un proyecto público, su realización estará justificada económicamente si su beneficio social marginal es superior a su coste directo en recursos multiplicado por el CMFP, en la medida que este es el precio social de cada euro recaudado con impuestos. 

En términos marginales, cuando el incremento en la recaudación es infinitesimal (o muy pequeño), la expresión ( 17) puede afinarse. tal que:

CMFP. = avt:(p,)
P' 8T(p',p,Y) (26)

donde VE es el coste de bienestar marginal medido a través del concepto expuesto de variación equivalente y que toma como referenc ia el precio previo a la aplicación del impuesto (p'), siendo Y la renta del individuo i. Alternativamente, se puede expresar el CMFP tomando como referencia el precio posterior al establecimiento del impuesto:

CMFP. =
avt:(p)
p iJT(p',p,Y) (27)

Obviamente los valores de (26) y (27) serán distintos, aunque. como señala González-Páramo (2003), ambas medidas son válidas, siempre y cuando sean consistentes, respectivamente, con los vectores de precios utilizados para evaluar los beneficios marginales del gasto público financiados con el aumento de recaudación evaluado. No obstante, hay que advertir que, para el propósito de la TlO, el coste marginal en bienestar derivado de financiar el objetivo de gasto público (en bienes públicos separables) con un sistema impositivo óptimo exige igualar los CMFP de todos los tipos de gravamen evaluados en el vector de precios posterior a su aplicación (p' = p'), lo que implica medir el CMFP a través de la expresión (27).


\subsection{Evasión fiscal}
La teoría de la evasión fiscal óptima desarrollada por ALLINGHAM y SANDMO (1972) es una contribución importante en el campo de la economía de la imposición. Esta teoría, cuyo antecedente es la teoría del comportamiento criminal de BECKER (1968), busca explicar cómo los individuos toman decisiones sobre si evadir o no impuestos y cómo el diseño del sistema tributario afecta sus incentivos para hacerlo. 

En concreto, la teoría se basa en la idea de que los individuos se enfrentan al dilema entre pagar sus impuestos legalmente. lo que reduce su ingreso disponible. y evadir impuestos. lo que les permite conservar más ingresos. pero con un riesgo de ser inspeccionados y. consecuentemente. de tener que pagar una sanción. Por lo tanto, los individuos toman la decisión sobre la cantidad óptima de impuesto a evadir (y, complementariamente, a declarar) teniendo en cuenta tanto la preferencia del individuo por la renta disponible tras el pago del impuesto, como el riesgo de ser comprobado, con la consiguien1e regularización del impuesto correcto y el pago de la sanción correspondiente. 

En su formulación básica. el modelo de evasión fiscal óptima considera un individuo con una renta (procedente del trabajo) y = w • L, por la que debe pagar un impuesto T(Y) = t · y. donde t es el tipo de brravamen. tal que O < t < l. El contribuyente puede elegir declarar una cuantía de renta x, igual o inferior a la renta obtenida (X y). de manera que estaría evadiendo una cantidad de renta e = y - x. En la aplicación de este impuesto. la administración tributaria dedica recursos para la comprobación de las declaraciones de renta. de manera que existe una probabilidad p de que el individuo sea inspeccionado, pudiendo no serlo con la probabilidad complementaria 1 - p. En el caso de que la declaración resulte inspeccionada. se supone que la actuación de la administración es plenamente eficaz, llevando a la regularización de toda la diferencia entre la renta declarada y la ganada. aplicándose una sanción s, tal que s > t. que en el modelo original se detine como una proporción de la renta evadida. s • e. Alternativamente, algunas variantes posteriores del modelo. como la de Yitzhaki (1974), definen la sanción como como un porcentajea aplicado sobre la cuota impositiva evadida, tal que a · t · e, siendo a > l.

De acuerdo con estas definiciones. un contribuyente que considera la posibilidad ele incumplir sus obligaciones fiscales decide qué cantidad de renta le interesa evadir, resolviendo el siguiente problema de maxímización de su utilidad esperada: 

max E(U) = (1 - p) • U(w • L • (l - t) + t · e) + p • U (w • L • (l - t) - s · e) (5.0) e donde la función de utilidad U(-), con U' (-) > O; U"(·) < O, refl<::ia sus preferencias por la renta con aversión al riesgo. La solución de este problema muestra la cantidad óptima de renta e• a evadir, tal que O \$ e' < y. de acuerdo con la condición de primer orden:\footnote{%
    La condición de segundo orden que garantiza. dada la concavidad de U(-). el máximo valor de E(U) es: (1 - p) • U"(w • L • (1 - t) + t · e) · t2 + p • U"(w • L • (1 - t) - s · e) • (s - t)2 < O.
}
U1(w·L·( l-t)-s·e) ( 1 -p)·t
u,(w·L· ( l -t)+t·e) p·(s-t) (51 )

De la condición anterior pueden deducirse algunas relaciones de interés para el diseño óptimo de la política de administración tributaria: a) El porcentaje de sanción actúa de forma disuasoria, de manera que al aumentar este se reduce la cantidad evadida óptima (oe/as < O). b) La probabilidad de ser inspeccionado también tiene carácter disuasorio. de manera que cuando esta aumenta la cantidad evadida óptima disminuye (oe/ éJp < O). c) Al incrementarse la renta salarial, la cantidad evadida óptima aumentará (éJe/éJy > O).
Esto se explica por la m��jor posición económica del contribuyente para afrontar actividades de riesgo, aunque el mismo resultado se obtiene también en el caso de que la aversión absoluta al riesgo sea decreciente.
d) El aumento del tipo impositivo tienen un efecto ambiguo sobre la evasión (oe/éJt �� O). 
Esto es debido a que, por un lado, un incremento del tipo impositivo provoca un "efecto renta·', ya que al reducirse la renta disponible de los contribuyentes que cumplan con sus obligaciones fiscales. si la aversión absoluta al riesgo es decreciente. estos mostrarán una menor disposición a evadir renta. En cambio, por otro lado. un aumento en el tipo impositivo origina un ·'efecto sustitución'· que hace más atractiva la evasión, pues con un tipo de gravamen más elevado el ahotTo fiscal por cada curo de renta no declarado será mayor, sin que aumente el coste unitario de la sanción en caso de inspección, si esta se calcula sobre la renta evadida. en lugar de sobre lacuota no satisfecha. En consecuencia, el impacto final dependerá de cuál de los dos efoctos  contrapuestos predomine.

Más adelante, YITZHAKI (1974) demuestra que si se establece una sanción (a > 1) sobre la cuota no satisfecha con la evasión. tal que. en caso de inspección, la renta disponible del contribuyente será, w • L - t · (w • L - e) - a · t • e = (1 - t) · w · L - (a - 1) • t • e, la condición de primer orden recogida en la expresión (51) pasa a ser:
u,((1-t)-w·L-(u-1 )·t·e) ( 1 -p)
u , (w·I,·(1-t)+c·e) = p·(o--1) (52)

Con este cambio en el mecanismo sancionador, la relación marginal con la que el contribuyente valora comparativamente la utilidad de la renta disponible según sea o no inspeccionado (primer término de la expresión) es ahora independiente del tipo impositivo (segundo término de la expresión). Por tanto, aunque pueda parecer contraintuitivo, cualquier aumento del tipo impositivo hace disminuir el importe óptimo de renta evadida. Esto es debido a que. al establecerse la sanción en función de la cuota no ingresada. la relación entre la evasión y el tipo impositivo solo presenta el efecto renta apuntado (ae/éJt < O).

\subsection{Teoría de la Administración Tributaria}
Recientemente, KEEN y SLEMROD (2017) han propuesto un marco de análisis para analizar el comportamiento óptimo de la administración tributaria, de forma integrada con los postulados de la TIO. El elemento fundamental es la incorporación del concepto denominado "elasticidad recaudatoria de la aplicación de la normativa impositiva" (e1?forceme111 elasticity of1ax revenue), con el que se estaría recogiendo la respuesta concluctual a las intervenciones de la administración tributa1ia. Esta elasticidad se diferencia d�� la elasticidad de la renta grava ble (elaslicity of tuxable income) que se limita a recoger las respuestas de la renta declarada a los cambios en los tipos marginales impositivos. incluidas las respuestas de evasión y elusión fiscales.\footnote{%
    En Saez, Slemrod )' Giertz (2012) se ofrece una completa revisión del concepto de elasticidad de la rl'.nta gravable.
} Este nuevo concepto de elasticidad permite la caracterización del equilibrio óptimo entre las medidas de diseño impositivo y las medidas de administración tributaria. así como el establecimiento de la brecha óptima de cumplimiento (optima/ co111p/iance gap).

En este modelo de administración tributaria óptima se incorpora la actividad de administración tributaria en términos de intensidad (a) al problema de elección del tipo impositivo óptimo t. baj o el supuesto realista de que conseguir bases impositivas que gravar (renta declarada) es una actividad costosa. El coste de esta actividad se explica porque los contribuyentes se comportan de acuerdo con una estrategia determinante de un nivel óptimo ele evasión fiscal, como hemos visto. y la administración tributaria 1iene que dedicar recursos para aumentar su capacidad ele control tributario. aumentado la probabilidad de detección ele la evasión. Además. en el modelo de Keen y Slemrod (2017). la fij ación del tipo impositivo óptimo tiene en cuenta el exceso ele gravamen a través de la elasticidad de la renta gravable, de acuerdo con la propuesta de Feldstein ( 1999). De este modo, la renta disponible de un individuo sometido a un impuesto t se define como,
w • L - t · (w • L - e) - c(e, a ) (53)

donde la renta efectivamente gravada es z = w • L - e, y e(-) es el coste de cumplimiento pr ivado al que se enfrenta el contribuyente a la hora de planificar su estrategia óptima de evasión/elusión. siendo argumentos de este coste el nivel de evasión e. y la intensidad de la actuación de la administración tributaria (a:). 14 

De acuerdo con estas de finiciones, el modelo plantea el siguiente problema de ma:ximización del bienestar social (W) para un individuo rep(esentativo (lo que supone no incorporar ninguna clase de condicionantes distTibuti vos):
max W(t, a) = w • L - t · (w • L - e) - c(e, a) - ijJ(L ) + V (g )
t,a
(54)

donde ijJ ( L) representa la valoración que el individuo hace del coste de oportunidad del ocio (medido en horas de trabajo ofertadas) y V (g) la valoración correspondiente al gasto público, que se financia con un impuesto sobre la renta salarial con tipo proporcional t. Para aplicar este impuesto. la administración tributaria incurre en un coste de administración a(a), creciente con la intensidad de intervenc ión (aa > O). En consecuencia, la restricción presupuestaria a la que se enfrenta el gobierno es:
g + a(a) = t · (w · L - e) (SS)

La resolución del programa for mado por (54) y (55) proporciona dos condiciones de primer orden. La primera muestra el impacto que tiene en el bienestar un cambio en el tipo impositivo, asumiendo el carácter distorsionante del impuesto sobre la renta.

t V'(·)-1 l - = -- · --- 1-1 V '(-J t(z,1-t)

donde c(z, 1 - t) es la elasticidad de la renta gravable al tipo marginal neto.\footnote{%
    El tipo marginal neto se interpreta como euro neto tra�� aplicar el tipo marginal.
} De acuerdo con esta condición. vemos que cuanto mayor sea el valor de esta elasticidad. mayor será el exceso de gravamen y, consecuentemente, más bajo debería ser el tipo impositivo óptimo.

La segunda condición es determinante del nivel óptimo del esfuerzo (tamaño) de la administración tributaria:

(57)

La interpretación intuitiva de esta condición es que la recaudación adicional que se puede obtener aumentando la intensidad de actuación por parte la administración tributaria debe equipararse a los costes adicionales de cumplimiento y administración asociados a la misma. siendo los segundos más relevantes que los primeros, pues los costes de administración se pagan con impuestos distorsionantes.

A partir de la condición (57 ), Keen y Slem rod (20 17) obtienen la mencionada '·elasticidad recaudatoria de la aplicación de la normativa impositiva·'. E(z,a),
a( C(t 1 {I ) ( ) v1(L·za-aa) ª E z, a = t·z
(58)

donde el numerador representa una combinación positiva de los costes de cum plimiento y de administración, con un menor peso de los primeros por el motivo indicado, mientras que el denominador recoge la recaudación obtenida al aplicar el tipo im positivo a la renta efectivamentegravada.

Ambas condiciones de optimalidad permiten obtener una serie de reglas de diseño óptimo tanto de las intervenciones de la administración tributaria como de política im positiva, donde se pueden destacar las tres siguientes:

a) Para una aplicación impositiva óptima por parte de la administración tributaria, es condición suficiente. aunque no necesaria, que el tipo impositivo (t) y la intensidad de actuación de la administración tributaria ( a) estén relacionados tal que, al aumentar la segunda. aumente el tipo (otjaa > O), disminuyendo la renta evadida (ea < O).

b) En el margen, es m��jor obtener ingresos impositivos adicionales. netos de costes de administración. incrementando el esfuerzo de aplicación de la normativa por parte de la administración tributaria. que aumentando el tipo impositivo. Esto se explica por los costes de eficiencia que apareja esta segunda opción.

e) El nivel óptimo del esfuerzo de administración tributaria se ve condicionado por las posibles relacione:. de complementariedad o sustituibilidad entre la evasión y la elusión que presente eldiseño impositivo.\footnote{%
    Esta regla se obtiene i ncorporando un importe q de la renta salarial ganada (wL) que no resu lta gravado por el impuesto, al estar afectado por prácticas de elusión fiscal. Estas actuacione�� implican costes de asesoramiento y transformación de rentas para el contribuyente, que pueden ser modelizados de acuerdo con la propuesta de Slemrod (2001).
}














































\clearpage
\section{Conclusión} 


\subsection{Recapitulación}


\subsection{Extensiones y relación con otras partes del Temario}



\subsection{Opinión}

\subsubsection{Idea final (Salida o cierra)}

\clearpage
\pagenumbering{gobble}
\MyBibliography



\MyBibItem{DAWID y GATTI}{Chapter 2 - Agent-Based Macroeconomics}{2018}{https://doi.org/10.1016/bs.hescom.2018.02.006}


% ANEXOS
\clearpage










\end{document}